% Source: Evans, "Partial Differential Equations" (2nd ed., AMS Graduate Studies in Mathematics vol. 19, 2010),
% Chapter 5 "Sobolev Spaces", PDF pages 245-298 of MathBooks/Evans Partial Differential Equations(1).pdf.
% Transcribed page-by-page by vision subagents, then merged and restructured into
% leanblueprint conventions (semantic \label, bracketed titles, \uses dependency edges) below.
%
% NOTE ON GAPS: as in chapter1.tex/chapter2.tex/chapter3.tex/chapter4.tex, a number of source
% pages were returned by the vision transcription subagent as refusal placeholders rather than
% actual page content (PDF pages 247, 259, 262, 265, 269, 277, 280, 284, 289, 295). As in
% chapter4.tex, every such gap was closed by directly re-reading the corresponding pages of the
% source PDF with a vision-capable Read call during this restructuring pass, so no content below
% is invented or reconstructed purely from outside knowledge of the book -- it is a faithful
% transcription of the actual page images. Running page headers (the small "5.X SECTION NAME"
% banner repeated at the top of every printed page) are omitted where they merely repeat an
% already-transcribed section title; only genuine inline section/subsection headings are kept.


This chapter mostly develops the theory of \textit{Sobolev spaces}, which turn out often to be the proper setting in which to apply ideas of functional analysis to glean information concerning partial differential equations. The following material is often subtle, and will seem largely unmotivated, but ultimately will prove extremely useful.

Since we have in mind eventual applications to rather wide classes of partial differential equations, it is worth sketching out here our overall point of view. Our intention, broadly put, will be later to take various specific PDE and to recast them abstractly as operators acting on appropriate linear spaces. We can symbolically write this as
$$A : X \to Y,$$
where the operator $A$ encodes the structure of the partial differential equations, including possibly boundary conditions, etc., and $X, Y$ are spaces of functions. The great advantage is that once our PDE problem has been suitably interpreted in this form, we can often employ the general and elegant principles of functional analysis (Appendix D) to study the solvability of various equations involving $A$. We will later see that the really hard work is not so much the invocation of functional analysis, but rather finding the ``right'' spaces $X, Y$ and the ``right'' abstract operators $A$. Sobolev spaces are designed precisely to make all this work out properly, and so these are usually the proper choices for $X, Y$.

We will utilize Sobolev spaces for studying linear elliptic, parabolic and hyperbolic PDE in Chapters 6-7, and for studying nonlinear elliptic and parabolic equations in Chapters 8-9.

The reader may wish to look over some of the terminology for functional analysis in Appendix D before going further.

\section{Hölder Spaces}
\label{sec:holder-spaces}

Before turning to Sobolev spaces, we first discuss the simpler \textit{Hölder spaces}.

Assume $U \subset \R^n$ is open and $0 < \gamma \le 1$. We have previously considered the class of Lipschitz continuous functions $u : U \to \R$, which by definition satisfy the estimate
$$(1) \quad |u(x) - u(y)| \le C|x - y| \quad (x,y \in U)$$
for some constant $C$. Now (1) of course implies $u$ is continuous, and more importantly provides a uniform modulus of continuity. It turns out to be useful to consider also functions $u$ satisfying a variant of (1), namely
$$(2) \quad |u(x) - u(y)| \le C|x - y|^\gamma \quad (x,y \in U)$$
for some constant $C$. Such a function is said to be \textit{Hölder continuous with exponent} $\gamma$.

\begin{definition}[Hölder seminorm and sup norm]
\label{def:holder-seminorm-sup-norm}
\dcref{ch5:5.1-def-1}
\group{Hölder Spaces}
(i) If $u : U \to \R$ is bounded and continuous, we write
$$\|u\|_{C(U)} := \sup_{x \in U} |u(x)|.$$
(ii) The $\gamma^{\text{th}}$-Hölder seminorm of $u : U \to \R$ is
$$[u]_{C^{0,\gamma}(U)} := \sup_{\substack{x,y \in U \\ x \ne y}} \left\{ \frac{|u(x) - u(y)|}{|x - y|^\gamma} \right\},$$
and the $\gamma^{\text{th}}$-Hölder norm is
$$\|u\|_{C^{0,\gamma}(\bar U)} := \|u\|_{C(\bar U)} + [u]_{C^{0,\gamma}(\bar U)}.$$
\end{definition}

\begin{definition}[Hölder space $C^{k,\gamma}(\bar U)$]
\label{def:holder-space}
\dcref{ch5:5.1-def-2}
\group{Hölder Spaces}
\uses{def:holder-seminorm-sup-norm}
The Hölder space
$$C^{k,\gamma}(\bar U)$$
consists of all functions $u \in C^k(\bar U)$ for which the norm
$$(3) \quad \|u\|_{C^{k,\gamma}(\bar U)} := \sum_{|\alpha| \le k} \|D^\alpha u\|_{C(\bar U)} + \sum_{|\alpha| = k} [D^\alpha u]_{C^{0,\gamma}(\bar U)}$$
is finite.
\end{definition}

So the space $C^{k,\gamma}(\bar U)$ consists of those functions $u$ that are $k$-times continuously differentiable and whose $k^{\text{th}}$-partial derivatives are Hölder continuous with exponent $\gamma$. Such functions are well-behaved, and furthermore the space $C^{k,\gamma}(\bar U)$ itself possesses a good mathematical structure:

\begin{theorem}[Hölder spaces as function spaces]
\label{thm:holder-space-banach}
\dcref{ch5:5.1-thm-1}
\group{Hölder Spaces}
\uses{def:holder-space}
The space of functions $C^{k,\gamma}(\bar U)$ is a Banach space.
\end{theorem}

The proof is left as an exercise (Problem 1), but let us pause here to make clear what is being asserted. Recall from §D.1 that if $X$ denotes a real linear space, then a mapping $\|\ \| : X \to [0,\infty)$ is called a \textit{norm} provided
\begin{itemize}
\item[(i)] $\|u+v\| \le \|u\| + \|v\|$ for all $u,v \in X$,
\item[(ii)] $\|\lambda u\| = |\lambda|\|u\|$ for all $u \in X$, $\lambda \in \R$,
\item[(iii)] $\|u\| = 0$ if and only if $u = 0$.
\end{itemize}
A norm provides us with a notion of convergence: we say a sequence $\{u_k\}_{k=1}^\infty \subset X$ converges to $u \in X$, written $u_k \to u$, if $\lim_{k\to\infty} \|u_k - u\| = 0$. A \textit{Banach space} is then a normed linear space that is \textit{complete}, that is, within which each Cauchy sequence converges.

So in Theorem \ref{thm:holder-space-banach} we are stating that if we take on the linear space $C^{k,\gamma}(\bar U)$ the norm $\|\cdot\| = \|\cdot\|_{C^{k,\gamma}(\bar U)}$, defined by (3), then $\|\cdot\|$ verifies properties (i)--(iii) above, and in addition each Cauchy sequence converges.

\section{Sobolev Spaces}
\label{sec:sobolev-spaces}

The Hölder spaces introduced in \cref{sec:holder-spaces} are unfortunately not often suitable settings for elementary PDE theory, as we usually cannot make good enough analytic estimates to demonstrate that the solutions we construct actually belong to such spaces. What are needed rather are some other kinds of spaces, containing less smooth functions. In practice we must strike a balance, by designing spaces comprising functions which have some, but not too great, smoothness properties.

\subsection{Weak derivatives}

We start off by substantially weakening the notion of partial derivatives.

\textbf{Notation.} Let $C_c^\infty(U)$ denote the space of infinitely differentiable functions $\phi: U \to \R$, with compact support in $U$. We will call a function $\phi$ belonging to $C_c^\infty(U)$ a \textit{test function}.

\textbf{Motivation for definition of weak derivative.} Assume we are given a function $u \in C^1(U)$. Then if $\phi \in C_c^\infty(U)$, we see from the integration by parts formula that
$$(1) \quad \int_U u\phi_{x_i}\,dx = -\int_U u_{x_i}\phi\,dx \quad (i=1,\ldots,n).$$
There are no boundary terms, since $\phi$ has compact support in $U$ and thus vanishes near $\partial U$. More generally now, if $k$ is a positive integer, $u \in C^k(U)$, and $\alpha = (\alpha_1,\ldots,\alpha_n)$ is a multiindex of order $|\alpha| = \alpha_1 + \cdots + \alpha_n = k$, then
$$(2) \quad \int_U uD^\alpha\phi\,dx = (-1)^{|\alpha|}\int_U D^\alpha u\,\phi\,dx.$$
This equality holds since
$$D^\alpha\phi = \frac{\partial^{\alpha_1}}{\partial x_1^{\alpha_1}} \cdots \frac{\partial^{\alpha_n}}{\partial x_n^{\alpha_n}}\phi$$
and we can apply formula (1) $|\alpha|$ times.

We next examine formula (2), valid for $u \in C^k(U)$, and ask whether some variant of it might still be true even if $u$ is not $k$ times continuously differentiable. Now the left hand side of (2) makes sense if $u$ is only locally summable: the problem is rather that if $u$ is not $C^k$, then the expression ``$D^\alpha u$'' on the right hand side of (2) has no obvious meaning. We resolve this difficulty by asking if there exists a locally summable function $v$ for which formula (2) is valid, with $v$ replacing $D^\alpha u$:

\begin{definition}[Weak partial derivative]
\label{def:weak-partial-derivative}
\dcref{ch5:5.2-def-1}
\group{Sobolev Spaces}
Suppose $u, v \in L^1_{\text{loc}}(U)$, and $\alpha$ is a multiindex. We say that $v$ is the $\alpha^{\text{th}}$-weak partial derivative of $u$, written
$$D^\alpha u = v,$$
provided
$$(3) \quad \int_U u D^\alpha \phi \, dx = (-1)^{|\alpha|} \int_U v\phi \, dx$$
for all test functions $\phi \in C_c^\infty(U)$.
\end{definition}

In other words, if we are given $u$ and if there happens to exist a function $v$ which verifies (3) for all $\phi$, we say that $D^\alpha u = v$ in the weak sense. If there does not exist such a function $v$, then $u$ does not possess a weak $\alpha^{\text{th}}$-partial derivative.

\begin{lemma}[Uniqueness of weak derivatives]
\label{lem:weak-derivative-uniqueness}
\dcref{ch5:5.2-lem-1}
\group{Sobolev Spaces}
\uses{def:weak-partial-derivative}
A weak $\alpha^{\text{th}}$-partial derivative of $u$, if it exists, is uniquely defined up to a set of measure zero.
\end{lemma}

\begin{proof}
Assume that $v, \bar{v} \in L^1_{\text{loc}}(U)$ satisfy
$$\int_U u D^\alpha \phi \, dx = (-1)^{|\alpha|} \int_U v\phi \, dx = (-1)^{|\alpha|} \int_U \bar{v}\phi \, dx$$
for all $\phi \in C_c^\infty(U)$. Then
$$(4) \quad \int_U (v - \bar{v})\phi \, dx = 0$$
for all $\phi \in C_c^\infty(U)$; whence $v - \bar{v} = 0$ a.e.
\end{proof}

\begin{example}[A weak derivative that exists]
\label{ex:weak-derivative-piecewise-linear-example}
\dcref{ch5:5.2-ex-1}
\group{Sobolev Spaces}
\uses{def:weak-partial-derivative}
Let $n = 1$, $U = (0,2)$, and
$$u(x) = \begin{cases}
x & \text{if } 0 < x \le 1 \\
1 & \text{if } 1 \le x < 2.
\end{cases}$$
Define
$$v(x) = \begin{cases}
1 & \text{if } 0 < x \le 1 \\
0 & \text{if } 1 < x < 2.
\end{cases}$$
Let us show $u' = v$ in the weak sense. To see this, choose any $\phi \in C_c^\infty(U)$. We must demonstrate
$$\int_0^2 u\phi' \, dx = -\int_0^2 v\phi \, dx.$$
But we easily calculate
$$\int_0^2 u\phi' \, dx = \int_0^1 x\phi' \, dx + \int_1^2 \phi' \, dx$$
$$= -\int_0^1 \phi \, dx + \phi(1) - \phi(1) = -\int_0^2 v\phi \, dx,$$
as required.
\end{example}

\begin{example}[A weak derivative that fails to exist]
\label{ex:weak-derivative-nonexistence-example}
\dcref{ch5:5.2-ex-2}
\group{Sobolev Spaces}
\uses{def:weak-partial-derivative}
Let $n = 1$, $U = (0,2)$, and
$$u(x) = \begin{cases} x & \text{if } 0 < x \le 1 \\ 2 & \text{if } 1 < x < 2. \end{cases}$$
We assert $u'$ does not exist in the weak sense. To check this, we must show there does not exist any function $v \in L^1_{\text{loc}}(U)$ satisfying
$$(5) \quad \int_0^2 u\phi' \, dx = -\int_0^2 v\phi \, dx$$
for all $\phi \in C_c^\infty(U)$. Suppose, to the contrary, (5) were valid for some $v$ and all $\phi$. Then
$$(6) \quad -\int_0^2 v\phi \, dx = \int_0^2 u\phi' \, dx = \int_0^1 x\phi' \, dx + 2\int_1^2 \phi' \, dx = -\int_0^1 \phi \, dx - \phi(1).$$
Choose a sequence $\{\phi_m\}_{m=1}^\infty$ of smooth functions satisfying
$$0 \le \phi_m \le 1, \quad \phi_m(1) = 1, \quad \phi_m(x) \to 0 \text{ for all } x \ne 1.$$
Replacing $\phi$ by $\phi_m$ in (6) and sending $m \to \infty$, we discover
$$1 = \lim_{m \to \infty} \phi_m(1) = \lim_{m \to \infty} \left[\int_0^2 v\phi_m \, dx - \int_0^1 \phi_m \, dx\right] = 0,$$
a contradiction.
\end{example}

More sophisticated examples appear in the next section.

\subsection{Definition of Sobolev spaces}

Fix $1 \le p \le \infty$ and let $k$ be a nonnegative integer. We define now certain function spaces, whose members have weak derivatives of various orders lying in various $L^p$ spaces.

\begin{definition}[Sobolev space $W^{k,p}(U)$]
\label{def:sobolev-space-wkp}
\dcref{ch5:5.2-def-2}
\group{Sobolev Spaces}
\uses{def:weak-partial-derivative}
The Sobolev space
$$W^{k,p}(U)$$
consists of all locally summable functions $u : U \to \R$ such that for each multiindex $\alpha$ with $|\alpha| \le k$, $D^\alpha u$ exists in the weak sense and belongs to $L^p(U)$.
\end{definition}

\begin{remark}[Notation $H^k$ and identification a.e.]
\label{rem:hk-notation-identification}
\dcref{ch5:5.2-rem-1}
\group{Sobolev Spaces}
\uses{def:sobolev-space-wkp}
(i) If $p = 2$, we usually write
$$H^k(U) = W^{k,2}(U) \quad (k = 0, 1, \ldots).$$
The letter $H$ is used, since---as we will see---$H^k(U)$ is a Hilbert space. Note that $H^0(U) = L^2(U)$.

(ii) We henceforth identify functions in $W^{k,p}(U)$ which agree a.e.
\end{remark}

\begin{definition}[Norm on $W^{k,p}(U)$]
\label{def:sobolev-norm}
\dcref{ch5:5.2-def-3}
\group{Sobolev Spaces}
\uses{def:sobolev-space-wkp}
If $u \in W^{k,p}(U)$, we define its norm to be
$$\|u\|_{W^{k,p}(U)} := \begin{cases}
\left(\sum_{|\alpha| \le k} \int_U |D^\alpha u|^p\,dx\right)^{1/p} & (1 \le p < \infty) \\
\sum_{|\alpha| \le k} \text{ess sup}_U |D^\alpha u| & (p = \infty).
\end{cases}$$
\end{definition}

\begin{definition}[Convergence in $W^{k,p}(U)$ and $W^{k,p}_{\text{loc}}(U)$]
\label{def:sobolev-convergence}
\dcref{ch5:5.2-def-4}
\group{Sobolev Spaces}
\uses{def:sobolev-norm}
(i) Let $\{u_m\}_{m=1}^\infty$, $u \in W^{k,p}(U)$. We say $u_m$ converges to $u$ in $W^{k,p}(U)$, written
$$u_m \to u \quad \text{in } W^{k,p}(U),$$
provided
$$\lim_{m \to \infty} \|u_m - u\|_{W^{k,p}(U)} = 0.$$

(ii) We write
$$u_m \to u \quad \text{in } W^{k,p}_{\text{loc}}(U),$$
to mean
$$u_m \to u \quad \text{in } W^{k,p}(V)$$
for each $V \CC U$.
\end{definition}

\begin{definition}[The space $W_0^{k,p}(U)$]
\label{def:sobolev-space-zero-boundary}
\dcref{ch5:5.2-def-5}
\group{Sobolev Spaces}
\uses{def:sobolev-convergence}
We denote by
$$W_0^{k,p}(U)$$
the closure of $C_c^\infty(U)$ in $W^{k,p}(U)$.
\end{definition}

Thus $u \in W_0^{k,p}(U)$ if and only if there exist functions $u_m \in C_c^\infty(U)$ such that $u_m \to u$ in $W^{k,p}(U)$. We interpret $W_0^{k,p}(U)$ as comprising those functions $u \in W^{k,p}(U)$ such that
$$\text{``}D^\alpha u = 0 \text{ on } \partial U \text{''} \quad \text{for all } |\alpha| \le k - 1.$$
This will all be made clearer with the discussion of traces in \cref{sec:sobolev-traces}.

\textbf{Notation.} It is customary to write
$$H^k_0(U) = W^{k,2}_0(U).$$

We will see in the exercises that if $n = 1$ and $U$ is an open interval in $\R^1$, then $u \in W^{1,p}(U)$ if and only if $u$ equals a.e. an absolutely continuous function whose ordinary derivative (which exists a.e.) belongs to $L^p(U)$. Such a simple characterization is however only available for $n = 1$. In general a function can belong to a Sobolev space, and yet be discontinuous and/or unbounded.

\begin{example}[A power singularity in a Sobolev space]
\label{ex:power-singularity-sobolev-membership}
\dcref{ch5:5.2-ex-3}
\group{Sobolev Spaces}
\uses{def:sobolev-space-wkp,def:weak-partial-derivative}
Take $U = B^0(0,1)$, the open unit ball in $\R^n$, and
$$u(x) = |x|^{-\alpha} \quad (x \in U, x \ne 0).$$
For which values of $\alpha > 0, n, p$ does $u$ belong to $W^{1,p}(U)$? To answer, note first $u$ is smooth away from $0$, with
$$u_{x_i}(x) = \frac{-\alpha x_i}{|x|^{\alpha+2}} \quad (x \ne 0),$$
and so
$$|Du(x)| = \frac{\alpha}{|x|^{\alpha+1}} \quad (x \ne 0).$$
Let $\phi \in C_c^\infty(U)$ and fix $\varepsilon > 0$. Then
$$\int_{U - B(0,\varepsilon)} u\phi_{x_i}\,dx = -\int_{U - B(0,\varepsilon)} u_{x_i}\phi\,dx + \int_{\partial B(0,\varepsilon)} u\phi\nu^i\,dS,$$
$\nu = (\nu^1, \ldots, \nu^n)$ denoting the inward pointing normal on $\partial B(0,\varepsilon)$. Now if $\alpha + 1 < n$, $|Du(x)| \in L^1(U)$. In this case
$$\left|\int_{\partial B(0,\varepsilon)} u\phi\nu^i\,dS\right| \le \|\phi\|_{L^\infty} \int_{\partial B(0,\varepsilon)} \varepsilon^{-\alpha}\,dS \le C\varepsilon^{n-1-\alpha} \to 0.$$
Thus
$$\int_U u\phi_{x_i}\,dx = -\int_U u_{x_i}\phi\,dx$$
for all $\phi \in C_c^\infty(U)$, provided $0 \le \alpha < n - 1$. Furthermore $|Du(x)| = \frac{\alpha}{|x|^{\alpha+1}} \in L^p(U)$ if and only if $(\alpha+1)p < n$. Consequently $u \in W^{1,p}(U)$ if and only if $\alpha < \frac{n-p}{p}$. In particular $u \notin W^{1,p}(U)$ for each $p \ge n$.
\end{example}

\begin{example}[A dense countable sum of singularities]
\label{ex:dense-singularities-unbounded-sobolev}
\dcref{ch5:5.2-ex-4}
\group{Sobolev Spaces}
\uses{ex:power-singularity-sobolev-membership}
Let $\{r_k\}_{k=1}^{\infty}$ be a countable, dense subset of $U = B^0(0,1)$. Write
$$u(x) = \sum_{k=1}^{\infty} \frac{1}{2^k} |x - r_k|^{-\alpha} \quad (x \in U).$$
Then $u \in W^{1,p}(U)$ if and only if $\alpha < \frac{n-p}{p}$. If $0 < \alpha < \frac{n-p}{p}$, we see that $u$ belongs to $W^{1,p}(U)$ and yet is unbounded on each open subset of $U$.
\end{example}

This last example illustrates a fundamental fact of life, that although a function $u$ belonging to a Sobolev space possesses certain smoothness properties, it can still be rather badly behaved in other ways.

\subsection{Elementary properties}

Next we verify certain properties of weak derivatives. Note very carefully that whereas these various rules are obviously true for smooth functions, functions in Sobolev space are not necessarily smooth: we must always rely solely upon the definition of weak derivatives.

\begin{theorem}[Properties of weak derivatives]
\label{thm:weak-derivative-properties}
\dcref{ch5:5.2-thm-1}
\group{Sobolev Spaces}
\uses{def:sobolev-space-wkp}
Assume $u, v \in W^{k,p}(U)$, $|\alpha| \le k$. Then

(i) $D^{\alpha}u \in W^{k-|\alpha|,p}(U)$ and $D^{\beta}(D^{\alpha}u) = D^{\alpha}(D^{\beta}u) = D^{\alpha+\beta}u$ for all multiindices $\alpha, \beta$ with $|\alpha| + |\beta| \le k$.

(ii) For each $\lambda, \mu \in \R$, $\lambda u + \mu v \in W^{k,p}(U)$ and $D^{\alpha}(\lambda u + \mu v) = \lambda D^{\alpha}u + \mu D^{\alpha}v$, $|\alpha| \le k$.

(iii) If $V$ is an open subset of $U$, then $u \in W^{k,p}(V)$.

(iv) If $\zeta \in C_c^{\infty}(U)$, then $\zeta u \in W^{k,p}(U)$ and
$$(7) \quad D^{\alpha}(\zeta u) = \sum_{\beta \le \alpha} \binom{\alpha}{\beta} D^{\beta}\zeta\, D^{\alpha-\beta}u \quad (\textit{Leibniz' formula}),$$
where $\binom{\alpha}{\beta} = \frac{\alpha!}{\beta!(\alpha-\beta)!}$.
\end{theorem}

\begin{proof}
1. To prove (i), first fix $\phi \in C_c^{\infty}(U)$. Then $D^{\beta}\phi \in C_c^{\infty}(U)$, and so
$$\int_U D^{\alpha}u\, D^{\beta}\phi \, dx = (-1)^{|\alpha|} \int_U u\, D^{\alpha+\beta}\phi \, dx$$
$$= (-1)^{|\alpha|}(-1)^{|\alpha+\beta|} \int_U D^{\alpha+\beta}u\,\phi \, dx$$
$$= (-1)^{|\beta|} \int_U D^{\alpha+\beta}u\,\phi \, dx.$$
Thus $D^\beta(D^\alpha u) = D^{\alpha+\beta} u$ in the weak sense.

2. Assertions (ii) and (iii) are easy, and the proofs are omitted.

3. We prove (7) by induction on $|\alpha|$. Suppose first $|\alpha| = 1$. Choose any $\varphi \in C_c^\infty(U)$. Then
$$\int_U \zeta u D^\alpha \varphi\,dx = \int_U u D^\alpha(\zeta\varphi) - u(D^\alpha \zeta)\varphi\,dx$$
$$= -\int_U (\zeta D^\alpha u + uD^\alpha\zeta)\varphi\,dx.$$
Thus $D^\alpha(\zeta u) = \zeta D^\alpha u + uD^\alpha\zeta$, as required.

Next assume $l < k$ and formula (7) is valid for all $|\alpha| \le l$ and all functions $\zeta$. Choose a multiindex $\alpha$ with $|\alpha| = l + 1$. Then $\alpha = \beta + \gamma$ for some $|\beta| = l$, $|\gamma| = 1$. Then for $\varphi$ as above,
$$\int_U \zeta u D^\alpha \varphi\,dx = \int_U \zeta u D^\beta(D^\gamma\varphi)\,dx$$
$$= (-1)^{|\beta|} \int_U \sum_{\sigma \le \beta} \binom{\beta}{\sigma} D^\sigma\zeta\, D^{\beta-\sigma}u\, D^\gamma\varphi\,dx$$
(by the induction assumption)
$$= (-1)^{|\beta|+|\gamma|} \int_U \sum_{\sigma \le \beta} \binom{\beta}{\sigma} D^\gamma(D^\sigma \zeta\, D^{\beta-\sigma}u)\varphi\,dx$$
(by the induction assumption again)
$$= (-1)^{|\alpha|} \int_U \sum_{\sigma \le \beta} \binom{\beta}{\sigma} [D^\sigma\zeta\, D^{\alpha-\rho}u + D^\gamma\zeta\, D^{\alpha-\sigma}u]\varphi\,dx$$
(where $\rho = \sigma + \gamma$)
$$= (-1)^{|\alpha|} \int_U \left[\sum_{\sigma \le \alpha} \binom{\alpha}{\sigma} D^\sigma\zeta\, D^{\alpha-\sigma}u\right] \varphi\,dx,$$
since
$$\binom{\beta}{\sigma - \gamma} + \binom{\beta}{\sigma} = \binom{\alpha}{\sigma}.$$

\begin{figure}[h]\centering
% [FIGURE: A small square symbol (QED mark)]
\end{figure}
\end{proof}

Not only do many of the usual rules of calculus apply to weak derivatives, but the Sobolev spaces themselves have a good mathematical structure:

\begin{theorem}[Sobolev spaces as function spaces]
\label{thm:sobolev-space-banach}
\dcref{ch5:5.2-thm-2}
\group{Sobolev Spaces}
\uses{def:sobolev-norm}
For each $k = 1, \ldots$ and $1 \le p \le \infty$, the Sobolev space $W^{k,p}(U)$ is a Banach space.
\end{theorem}

\begin{proof}
1. Let us first of all check that $\|u\|_{W^{k,p}(U)}$ is a norm. (See the discussion at the end of \cref{sec:holder-spaces}, or refer to §D.1, for definitions.) Clearly
$$\|\lambda u\|_{W^{k,p}(U)} = |\lambda| \|u\|_{W^{k,p}(U)},$$
and
$$\|u\|_{W^{k,p}(U)} = 0 \text{ if and only if } u = 0 \text{ a.e.}$$
Next assume $u, v \in W^{k,p}(U)$. Then if $1 \le p < \infty$, Minkowski's inequality (§B.2) implies
$$\|u + v\|_{W^{k,p}(U)} = \left(\sum_{|\alpha| \le k} \|D^\alpha u + D^\alpha v\|_{L^p(U)}^p\right)^{1/p}$$
$$\le \left(\sum_{|\alpha| \le k} (\|D^\alpha u\|_{L^p(U)} + \|D^\alpha v\|_{L^p(U)})^p\right)^{1/p}$$
$$\le \left(\sum_{|\alpha| \le k} \|D^\alpha u\|_{L^p(U)}^p\right)^{1/p} + \left(\sum_{|\alpha| \le k} \|D^\alpha v\|_{L^p(U)}^p\right)^{1/p}$$
$$= \|u\|_{W^{k,p}(U)} + \|v\|_{W^{k,p}(U)}.$$

2. It remains to show that $W^{k,p}(U)$ is complete. So assume $\{u_m\}_{m=1}^\infty$ is a Cauchy sequence in $W^{k,p}(U)$. Then for each $|\alpha| \le k$, $\{D^\alpha u_m\}_{m=1}^\infty$ is a Cauchy sequence in $L^p(U)$. Since $L^p(U)$ is complete, there exist functions $u_\alpha \in L^p(U)$ such that
$$D^\alpha u_m \to u_\alpha \text{ in } L^p(U)$$
for each $|\alpha| \le k$. In particular,
$$u_m \to u_{(0,\ldots,0)} =: u \text{ in } L^p(U).$$

3. We now claim
$$(8) \quad u \in W^{k,p}(U), \quad D^\alpha u = u_\alpha \quad (|\alpha| \le k).$$
To verify this assertion, fix $\phi \in C_c^\infty(U)$. Then
$$\int_U u D^\alpha \phi \, dx = \lim_{m \to \infty} \int_U u_m D^\alpha \phi \, dx$$
$$= \lim_{m \to \infty} (-1)^{|\alpha|} \int_U D^\alpha u_m \phi \, dx$$
$$= (-1)^{|\alpha|} \int_U u_\alpha \phi \, dx.$$
Thus (8) is valid. Since therefore $D^\alpha u_m \to D^\alpha u$ in $L^p(U)$ for all $|\alpha| \le k$, we see that $u_m \to u$ in $W^{k,p}(U)$, as required.
\end{proof}

\section{Approximation}
\label{sec:sobolev-approximation}

\subsection{Interior approximation by smooth functions}

It is awkward to return continually to the definition of weak derivatives. In order to study the deeper properties of Sobolev spaces, we therefore need to develop some systematic procedures for approximating a function in a Sobolev space by smooth functions. The method of mollifiers, developed in §C.4, provides the tool.

Fix a positive integer $k$ and $1 \le p < \infty$. Remember that $U_\varepsilon = \{x \in U \mid \dist(x, \partial U) > \varepsilon\}$.

\begin{theorem}[Local approximation by smooth functions]
\label{thm:local-smooth-approximation}
\dcref{ch5:5.3-thm-1}
\group{Sobolev Space Approximation}
\uses{def:sobolev-space-wkp,def:sobolev-convergence}
Assume $u \in W^{k,p}(U)$ for some $1 \le p < \infty$, and set
$$u^\varepsilon = \eta_\varepsilon * u \text{ in } U_\varepsilon.$$
Then
\begin{enumerate}
\item[(i)] $u^\varepsilon \in C^\infty(U_\varepsilon)$ for each $\varepsilon > 0$,
\end{enumerate}
and
\begin{enumerate}
\item[(ii)] $u^\varepsilon \to u$ in $W^{k,p}_{\text{loc}}(U)$, as $\varepsilon \to 0$.
\end{enumerate}
\end{theorem}

\begin{proof}
1. Assertion (i) is proved in §C.4.

2. We next claim that if $|\alpha| \le k$, then
$$(1) \quad D^\alpha u^\varepsilon = \eta_\varepsilon * D^\alpha u \text{ in } U_\varepsilon;$$
that is, the ordinary $\alpha^{\text{th}}$-partial derivative of the smooth function $u^\varepsilon$ is the $\varepsilon$-mollification of the $\alpha^{\text{th}}$-weak partial derivative of $u$. To confirm this, we compute for $x \in U_\varepsilon$
$$D^\alpha u^\varepsilon(x) = D^\alpha \int_U \eta_\varepsilon(x-y)u(y)\,dy$$
$$= \int_U D_x^\alpha \eta_\varepsilon(x-y)u(y)\,dy$$
$$= (-1)^{|\alpha|} \int_U D_y^\alpha \eta_\varepsilon(x-y)u(y)\,dy.$$
Now for fixed $x \in U_\varepsilon$ the function $\phi(y) := \eta_\varepsilon(x-y)$ belongs to $C_c^\infty(U)$. Consequently the definition of the $\alpha^{\text{th}}$-weak partial derivative implies:
$$\int_U D_y^\alpha \eta_\varepsilon(x-y)u(y)\,dy = (-1)^{|\alpha|} \int_U \eta_\varepsilon(x-y)D^\alpha u(y)\,dy.$$
Thus
$$D^\alpha u^\varepsilon(x) = (-1)^{|\alpha|+|\alpha|}\int_U \eta_\varepsilon(x-y)D^\alpha u(y)\,dy$$
$$= [\eta_\varepsilon * D^\alpha u](x).$$
This establishes (1).

3. Now choose an open set $V \CC U$. In view of (1) and §C.4, $D^\alpha u^\varepsilon \to D^\alpha u$ in $L^p(V)$ as $\varepsilon \to 0$, for each $|\alpha| \le k$. Consequently
$$\|u^\varepsilon - u\|_{W^{k,p}(V)}^p = \sum_{|\alpha|\le k} \|D^\alpha u^\varepsilon - D^\alpha u\|_{L^p(V)}^p \to 0$$
as $\varepsilon \to 0$. This proves assertion (ii).
\end{proof}

\subsection{Global approximation by smooth functions}

Next we show that we can find smooth functions which approximate in $W^{k,p}(U)$, and not just in $W^{k,p}_{\text{loc}}(U)$. Notice in the following that we make no assumptions about the smoothness of $\partial U$.

\begin{theorem}[Global approximation by smooth functions]
\label{thm:global-smooth-approximation-interior}
\dcref{ch5:5.3-thm-2}
\group{Sobolev Space Approximation}
\uses{def:sobolev-space-wkp,thm:weak-derivative-properties}
Assume $U$ is bounded, and suppose as well that $u \in W^{k,p}(U)$ for some $1 \le p < \infty$. Then there exist functions $u_m \in C^\infty(U) \cap W^{k,p}(U)$ such that
$$u_m \to u \quad \text{in } W^{k,p}(U).$$
\end{theorem}

\begin{remark}[Approximation not up to the boundary]
\label{rem:approximation-not-up-to-boundary}
\dcref{ch5:5.3-rem-1}
\group{Sobolev Space Approximation}
\uses{thm:global-smooth-approximation-interior}
Note carefully that we do not assert $u_m \in C^\infty(\bar U)$ (but see \cref{thm:global-smooth-approximation-boundary} below).
\end{remark}

\begin{proof}
1. We have $U = \bigcup_{i=1}^\infty U_i$, where
$$U_i := \{x \in U \mid \dist(x, \partial U) > 1/i\} \quad (i = 1, 2, \ldots).$$
Write $V_i := U_{i+3} - \overline{U_{i+1}}$.

Choose also any open set $V_0 \CC U$ so that $U = \bigcup_{i=0}^\infty V_i$. Now let $\{\zeta_i\}_{i=0}^\infty$ be a smooth partition of unity subordinate to the open sets $\{V_i\}_{i=0}^\infty$; that is, suppose
$$(2) \quad \begin{cases} 0 \le \zeta_i \le 1, \quad \zeta_i \in C_c^\infty(V_i) \\ \sum_{i=0}^\infty \zeta_i = 1 \quad \text{on } U. \end{cases}$$
Next, choose any function $u \in W^{k,p}(U)$. According to Theorem \ref{thm:weak-derivative-properties}(iv), $\zeta_i u \in W^{k,p}(U)$ and $\spt(\zeta_i u) \subset V_i$.

2. Fix $\delta > 0$. Choose then $\varepsilon_i > 0$ so small that $u^i := \eta_{\varepsilon_i} * (\zeta_i u)$ satisfies
$$(3) \quad \begin{cases}
\|u^i - \zeta_i u\|_{W^{k,p}(U)} \le \frac{\delta}{2^{i+1}} & (i = 0,1,\ldots) \\
\spt\, u^i \subset W_i & (i = 1,\ldots),
\end{cases}$$
for $W_i := U_{i+4} - \overline{U_i} \supset V_i$ $(i = 1,\ldots)$.

3. Write $v := \sum_{i=0}^\infty u^i$. This function belongs to $C^\infty(U)$, since for each open set $V \CC U$ there are at most finitely many nonzero terms in the sum. Since $u = \sum_{i=0}^\infty \zeta_i u$, we have for each $V \CC U$
$$\|v - u\|_{W^{k,p}(V)} \le \sum_{i=0}^\infty \|u^i - \zeta_i u\|_{W^{k,p}(U)}$$
$$\le \delta \sum_{i=0}^\infty \frac{1}{2^{i+1}} \text{ by (3)}$$
$$= \delta.$$
Take the supremum over sets $V \CC U$, to conclude $\|v - u\|_{W^{k,p}(U)} \le \delta$.
\end{proof}

\subsection{Global approximation by functions smooth up to the boundary}

We now ask when it is possible to approximate a given function $u \in W^{k,p}(U)$ by functions belonging to $C^\infty(\bar U)$, and not just $C^\infty(U)$. Such an approximation requires some condition to exclude $\partial U$ being wild geometrically.

\begin{theorem}[Global approximation by functions smooth up to the boundary]
\label{thm:global-smooth-approximation-boundary}
\dcref{ch5:5.3-thm-3}
\group{Sobolev Space Approximation}
\uses{def:sobolev-space-wkp,thm:local-smooth-approximation,thm:weak-derivative-properties}
Assume $U$ is bounded and $\partial U$ is $C^1$. Suppose $u \in W^{k,p}(U)$ for some $1 \le p < \infty$. Then there exist functions $u_m \in C^\infty(\bar U)$ such that
$$u_m \to u \quad \text{in } W^{k,p}(U).$$
\end{theorem}

\begin{proof}
1. Fix any point $x^0 \in \partial U$. As $\partial U$ is $C^1$, there exist, according to §C.1, a radius $r > 0$ and a $C^1$ function $\gamma : \R^{n-1} \to \R$ such that---upon relabeling the coordinate axes if necessary---we have
$$U \cap B(x^0, r) = \{x \in B(x^0, r) \mid x_n > \gamma(x_1, \ldots, x_{n-1})\}.$$
Set $V := U \cap B(x^0, r/2)$.

2. Define the shifted point
$$x^\varepsilon := x + \lambda \varepsilon e_n \quad (x \in V, \varepsilon > 0),$$
and observe that for some fixed, sufficiently large number $\lambda > 0$ the ball $B(x^\varepsilon, \varepsilon)$ lies in $U \cap B(x^0, r)$ for all $x \in V$ and all small $\varepsilon > 0$.

\begin{figure}[h]\centering
% [FIGURE: A coordinate diagram showing the domain U bounded below by the curve x_n = gamma(x_1,...,x_{n-1}), with the shifted ball B(x^varepsilon,varepsilon) sitting strictly inside U above a point x in the smaller half-ball V, and the larger ball B(x^0,r) centered at the boundary point x^0.]
\end{figure}

Now define $u_\varepsilon(x) := u(x^\varepsilon)$ ($x \in V$); this is the function $u$ translated a distance $\lambda\varepsilon$ in the $e_n$ direction. Next write $v^\varepsilon = \eta_\varepsilon * u_\varepsilon$. The idea is that we have moved up enough so that ``there is room to mollify within $U$''. Clearly $v^\varepsilon \in C^\infty(\bar V)$.

3. We now claim
$$(4) \quad v^\varepsilon \to u \quad \text{in } W^{k,p}(V).$$
To confirm this, take $\alpha$ to be any multiindex with $|\alpha| \le k$. Then
$$\|D^\alpha v^\varepsilon - D^\alpha u\|_{L^p(V)} \le \|D^\alpha v^\varepsilon - D^\alpha u_\varepsilon\|_{L^p(V)} + \|D^\alpha u_\varepsilon - D^\alpha u\|_{L^p(V)}.$$
The second term on the right hand side goes to zero with $\varepsilon$, since translation is continuous in the $L^p$-norm; and the first term also vanishes in the limit, by reasoning similar to that in the proof of Theorem \ref{thm:local-smooth-approximation}, for each $|\alpha| \le k$, and so (4) follows.

4. Using these calculations we readily check as well
$$(5) \quad \|v^\varepsilon\|_{W^{k,p}(V)} \le C\|u\|_{W^{k,p}(V)}$$
for some constant $C$ which does not depend on $u$.

5. Let us next consider the situation that $\partial U$ is not necessarily flat near $x^0$. Then utilizing the notation and terminology from §C.1, we can find a $C^1$ mapping $\Phi$, with inverse $\Psi$, such that $\Phi$ ``straightens out $\partial U$ near $x^0$''.

\begin{figure}[h]\centering
% [FIGURE: Two side-by-side disks labeled "x-coordinates" and "y-coordinates", connected by curved arrows labeled Phi (pointing right) and Psi (pointing left), captioned "Straightening out the boundary".]
\end{figure}

We write $y = \Phi(x)$, $x = \Psi(y)$, $u'(y) := u(\Psi(y))$. Choose a small ball $B$ as drawn before. Then utilizing steps 1--4 above we approximate $u'$ on $B^+ := B \cap \{y_n > 0\}$ by functions $v'^\varepsilon \in C^\infty(\bar B^+ \cap \bar B(x^0,r/2)$-type region$)$, such that $v'^\varepsilon \to u'$ in $W^{k,p}(B^+ \cap V')$, for $V'$ a smaller half-ball, with
$$\|v'^\varepsilon\|_{W^{k,p}(V')} \le C\|u'\|_{W^{k,p}(B^+)}.$$
Let $W := \Psi(V')$. Then converting back to the $x$-variables, we obtain a function $v^\varepsilon(x) := v'^\varepsilon(\Phi(x))$ approximating $u$ on $W$, with
$$(6) \quad \|v^\varepsilon\|_{W^{k,p}(W)} \le C\|u\|_{W^{k,p}(U)}.$$

6. Since $\partial U$ is compact, there exist finitely many points $x_i^0 \in \partial U$, radii $r_i > 0$, corresponding sets $V_i = U \cap B(x_i^0, r_i/2)$, and functions $v_i \in C^\infty(\bar V_i)$ ($i = 1,\ldots,N$) such that $\partial U \subset \bigcup_{i=1}^N B^0(x_i^0, r_i/2)$, and
$$\|v_i - u\|_{W^{k,p}(V_i)} \le \delta.$$
Take an open set $V_0 \CC U$ such that $U \subset \bigcup_{i=0}^N V_i$ and select, using Theorem \ref{thm:local-smooth-approximation}, a function $v_0 \in C^\infty(\bar V_0)$ satisfying
$$\|v_0 - u\|_{W^{k,p}(V_0)} \le \delta.$$

7. Now let $\{\zeta_i\}_{i=0}^N$ be a smooth partition of unity subordinate to the open sets $\{V_i\}_{i=0}^N$ in $U$. Define $v := \sum_{i=0}^N \zeta_i v_i$. Then clearly $v \in C^\infty(\bar U)$. In addition, since $u = \sum_{i=0}^N \zeta_i u$, we see using Theorem \ref{thm:weak-derivative-properties} that for each $|\alpha| \le k$:
$$\|D^\alpha v - D^\alpha u\|_{L^p(U)} \le \sum_{i=0}^N \|D^\alpha(\zeta_i v_i) - D^\alpha(\zeta_i u)\|_{L^p(V_i)}$$
$$\le C \sum_{i=0}^N \|v_i - u\|_{W^{k,p}(V_i)} = CN\delta,$$
according to the estimates above.
\end{proof}

\section{Extensions}
\label{sec:sobolev-extensions}

Our goal next is to extend functions in the Sobolev space $W^{1,p}(U)$ to become functions in the Sobolev space $W^{1,p}(\R^n)$. This can be subtle. Observe for instance that our extending $u \in W^{1,p}(U)$ to be zero in $\R^n - U$ will not in general work, as we may thereby create such a bad discontinuity along $\partial U$ that the extended function no longer has a weak first partial derivative. We must instead invent a way to extend $u$ which ``preserves the weak derivatives across $\partial U$''.

Suppose $1 \le p \le \infty$.

\begin{theorem}[Extension Theorem]
\label{thm:sobolev-extension-theorem}
\dcref{ch5:5.4-thm-1}
\group{Sobolev Extension}
\uses{def:sobolev-space-wkp,thm:global-smooth-approximation-boundary}
Assume $U$ is bounded and $\partial U$ is $C^1$. Select a bounded open set $V$ such that $U \CC V$. Then there exists a bounded linear operator
$$(1) \quad E: W^{1,p}(U) \to W^{1,p}(\R^n)$$
such that for each $u \in W^{1,p}(U)$:
\begin{itemize}
\item[(i)] $Eu = u$ a.e. in $U$,
\item[(ii)] $Eu$ has support within $V$,
\end{itemize}
and
\begin{itemize}
\item[(iii)] $\|Eu\|_{W^{1,p}(\R^n)} \le C\|u\|_{W^{1,p}(U)},$
\end{itemize}
the constant $C$ depending only on $p$, $U$, and $V$.
\end{theorem}

\begin{definition}[Extension of a Sobolev function]
\label{def:extension-of-sobolev-function}
\dcref{ch5:5.4-def-1}
\group{Sobolev Extension}
\uses{thm:sobolev-extension-theorem}
We call $Eu$ an extension of $u$ to $\R^n$.
\end{definition}

\begin{proof}
1. Fix $x^0 \in \partial U$ and suppose first
$$(2) \quad \partial U \text{ is flat near } x^0, \text{ lying in the plane } \{x_n = 0\}.$$

\begin{figure}[h]\centering
% [FIGURE: A coordinate system showing a shaded half-disk region representing the domain U above the horizontal axis, with a semicircular region B^- below the axis. The figure illustrates a half-ball at the boundary of a domain.]
\end{figure}

\textbf{A half-ball at the boundary}

Then we may assume there exists an open ball $B$, with center $x^0$ and radius $r$, such that
$$\begin{cases}
B^+ := B \cap \{x_n \ge 0\} \subset \bar{U} \\
B^- := B \cap \{x_n \le 0\} \subset \R^n - U.
\end{cases}$$

2. Temporarily suppose also $u \in C^\infty(\bar{U})$. We define then
$$(3) \quad \tilde{u}(x) := \begin{cases}
u(x) & \text{if } x \in B^+ \\
-3u(x_1, \ldots, x_{n-1}, -x_n) + 4u(x_1, \ldots, x_{n-1}, -\frac{x_n}{2}) & \text{if } x \in B^-.
\end{cases}$$
This is called a \textit{higher-order reflection} of $u$ from $B^+$ to $B^-$.

3. We claim
$$(4) \quad \tilde{u} \in C^1(B).$$
To check this, let us write $u^- := \tilde{u}|_{B^-}$, $u^+ := \tilde{u}|_{B^+}$. We demonstrate first
$$(5) \quad u^-_{x_n} = u^+_{x_n} \text{ on } \{x_n = 0\}.$$
Indeed according to (3),
$$\frac{\partial u^-}{\partial x_n}(x) = 3\frac{\partial u}{\partial x_n}(x_1, \ldots, x_{n-1}, -x_n) - 2\frac{\partial u}{\partial x_n}(x_1, \ldots, x_{n-1}, -\frac{x_n}{2})$$
and so
$$u^-_{x_n}|_{\{x_n=0\}} = u^+_{x_n}|_{\{x_n=0\}}.$$
This confirms (5). Now since $u^+ = u^-$ on $\{x_n = 0\}$, we see as well that
$$(6) \quad u^-_{x_i}|_{\{x_n=0\}} = u^+_{x_i}|_{\{x_n=0\}}$$
for $i = 1, \ldots, n-1$. But then (5) and (6) together imply
$$D^\alpha u^-|_{\{x_n=0\}} = D^\alpha u^+|_{\{x_n=0\}}$$
for each $|\alpha| \le 1$, and so (4) follows.

4. Using these calculations we readily check as well
$$(7) \quad \|\tilde u\|_{W^{1,p}(B)} \le C\|u\|_{W^{1,p}(B^+)}$$
for some constant $C$ which does not depend on $u$.

5. Let us next consider the situation that $\partial U$ is not necessarily flat near $x^0$. Then utilizing the notation and terminology from §C.1, we can find a $C^1$ mapping $\Phi$, with inverse $\Psi$, such that $\Phi$ ``straightens out $\partial U$ near $x^0$''.

\begin{figure}[h]\centering
% [FIGURE: Two side-by-side disks labeled "x-coordinates" and "y-coordinates", connected by curved arrows labeled Phi (pointing right) and Psi (pointing left), captioned "Straightening out the boundary".]
\end{figure}

\textbf{Straightening out the boundary}

We write $y = \Phi(x)$, $x = \Psi(y)$, $u'(y) := u(\Psi(y))$. Choose a small ball $B$ as drawn before. Then utilizing steps 1--3 above we extend $u'$ from $B^+$ to a function $\bar u'$ defined on all of $B$, such that $\bar u'$ is $C^1$ and we have the estimate
$$\|\bar u'\|_{W^{1,p}(B)} \le C\|u'\|_{W^{1,p}(B^+)}.$$
Let $W := \Psi(B)$. Then converting back to the $x$-variables, we obtain an extension $\bar u$ of $u$ to $W$, with
$$(8) \quad \|\bar u\|_{W^{1,p}(W)} \le C\|u\|_{W^{1,p}(U)}.$$

6. Since $\partial U$ is compact, there exist finitely many points $x_i^0 \in \partial U$, open sets $W_i$, and extensions $\bar u_i$ of $u$ to $W_i$ ($i = 1,\ldots,N$), as above, such that $\Gamma \subset \bigcup_{i=1}^N W_i$. Take $W_0 \CC U$ so that $U \subset \bigcup_{i=0}^N W_i$, and let $\{\zeta_i\}_{i=0}^N$ be an associated partition of unity. Write $\bar u := \sum_{i=0}^N \zeta_i \bar u_i$, where $\bar u_0 = u$. Then utilizing estimate (8) (with $u_i$ in place of $u$, $\bar u_i$ in place of $\bar u$) we obtain the bound
$$(9) \quad \|\bar u\|_{W^{1,p}(\R^n)} \le C\|u\|_{W^{1,p}(U)}$$
for some constant $C$, depending on $U,p,n$, etc., but not on $u$. Furthermore we can arrange for the support of $\bar u$ to lie within $V \supset\supset U$.

7. We henceforth write $Eu := \bar{u}$, and observe that the mapping $u \mapsto Eu$ is linear.

Recall that the construction so far assumed $u \in C^\infty(\bar{U})$. Suppose now $u \in W^{1,p}(U)$, and choose $u_m \in C^\infty(\bar{U})$ converging to $u$ in $W^{1,p}(U)$ (using Theorem \ref{thm:global-smooth-approximation-boundary}). Estimate (9) and the linearity of $E$ imply
$$\|Eu_m - Eu_l\|_{W^{1,p}(\R^n)} \le C\|u_m - u_l\|_{W^{1,p}(U)}.$$
Thus $\{Eu_m\}_{m=1}^\infty$ is a Cauchy sequence and so converges to $\bar{u} := Eu$. This extension, which does not depend on the particular choice of the approximating sequence $\{u_m\}_{m=1}^\infty$, satisfies the conclusions of the theorem.
\end{proof}

\begin{remark}[Extension in $W^{2,p}$; failure for $k>2$]
\label{rem:extension-higher-order-remark}
\dcref{ch5:5.4-rem-1}
\group{Sobolev Extension}
\uses{thm:sobolev-extension-theorem}
(i) Assume now that $\partial U$ is $C^2$. Then the extension operator $E$ constructed above is also a bounded linear operator from $W^{2,p}(U)$ to $W^{2,p}(\R^n)$. To see this, note first in steps 3, 4 of the proof that although $\bar{u}$ is not in general $C^2$, it does belong to $W^{2,p}(B)$. We also have the bound
$$\|\bar{u}\|_{W^{2,p}(B)} \le C\|u\|_{W^{2,p}(B_+)},$$
which follows from the definition (3). As before, we consequently derive the estimate
$$(10) \quad \|Eu\|_{W^{2,p}(\R^n)} \le C\|u\|_{W^{2,p}(U)},$$
provided $\partial U$ is $C^2$, the constants $C$ depending only on $U, V, n$ and $p$. We will need these observations later.

(ii) The above construction does \textit{not} provide us with an extension for the Sobolev spaces $W^{k,p}(U)$, if $k > 2$. This requires a more complicated higher-order reflection technique.
\end{remark}

\section{Traces}
\label{sec:sobolev-traces}

Next we discuss the possibility of assigning ``boundary values'' along $\partial U$ to a function $u \in W^{1,p}(U)$, assuming that $\partial U$ is $C^1$. Now if $u \in C(\bar{U})$, then clearly $u$ has values on $\partial U$ in the usual sense. The problem is that a typical function $u \in W^{1,p}(U)$ is not in general continuous and, even worse, is only defined a.e. in $U$. Since $\partial U$ has $n$-dimensional Lebesgue measure zero, there is no direct meaning we can give to the expression ``$u$ restricted to $\partial U$''. The notion of a \textit{trace operator} resolves this problem.

For this section we take $1 \le p < \infty$.

\begin{theorem}[Trace Theorem]
\label{thm:trace-theorem}
\dcref{ch5:5.5-thm-1}
\group{Sobolev Traces}
\uses{def:sobolev-space-wkp,thm:global-smooth-approximation-boundary}
Assume $U$ is bounded and $\partial U$ is $C^1$. Then there exists a bounded linear operator
$$T : W^{1,p}(U) \to L^p(\partial U)$$
such that
\begin{itemize}
\item[(i)] $Tu = u|_{\partial U}$ if $u \in W^{1,p}(U) \cap C(\overline{U})$
\end{itemize}
and
\begin{itemize}
\item[(ii)] $$\|Tu\|_{L^p(\partial U)} \le C\|u\|_{W^{1,p}(U)},$$
\end{itemize}
for each $u \in W^{1,p}(U)$, with the constant $C$ depending only on $p$ and $U$.
\end{theorem}

\begin{definition}[Trace of $u$ on $\partial U$]
\label{def:trace-operator}
\dcref{ch5:5.5-def-1}
\group{Sobolev Traces}
\uses{thm:trace-theorem}
We call $Tu$ the trace of $u$ on $\partial U$.
\end{definition}

\begin{proof}
1. Assume first $u \in C^1(\overline{U})$. As in the first part of the proof of Theorem \ref{thm:sobolev-extension-theorem}, let us also initially suppose $x^0 \in \partial U$ and $\partial U$ is flat near $x^0$, lying in the plane $\{x_n = 0\}$. Choose an open ball $B$ as in the previous proof and let $\bar{B}$ denote the concentric ball with radius $r/2$.

Select $\zeta \in C_c^\infty(B)$, with $\zeta \ge 0$ in $B$, $\zeta \equiv 1$ on $\bar{B}$. Denote by $\Gamma$ that portion of $\partial U$ within $\bar{B}$. Set $x' = (x_1, \ldots, x_{n-1}) \in \R^{n-1} = \{x_n = 0\}$.

Then
$$\int_\Gamma |u|^p\,dx' \le \int_{\{x_n=0\}} \zeta|u|^p\,dx' = -\int_{B^+} (\zeta|u|^p)_{x_n}\,dx$$
$$(1) \quad \begin{split}
& = -\int_{B^+} |u|^p\zeta_{x_n} + p|u|^{p-1}(\sgn u)u_{x_n}\zeta \, dx\\
& \le C\int_{B^+} |u|^p + |Du|^p \, dx,
\end{split}$$
where we employed Young's inequality, from §B.2.

2. If $x^0 \in \partial U$, but $\partial U$ is not flat near $x^0$, we as usual straighten out the boundary near $x^0$ to obtain the setting above. Applying estimate (1) and changing variables, we obtain the bound
$$\int_\Gamma |u|^p\,dS \le C\int_U |u|^p + |Du|^p \, dx,$$
where $\Gamma$ is some open subset of $\partial U$ containing $x^0$.

3. Since $\partial U$ is compact, there exist finitely many points $x_i^0 \in \partial U$ and open subsets $\Gamma_i \subset \partial U$ ($i = 1, \ldots, N$) such that $\partial U = \bigcup_{i=1}^N \Gamma_i$ and
$$\|u\|_{L^p(\Gamma_i)} \le C\|u\|_{W^{1,p}(U)} \quad (i = 1, \ldots, N).$$
Consequently, if we write
$$Tu := u|_{\partial U},$$
then
$$(2) \quad \|Tu\|_{L^p(\partial U)} \le C\|u\|_{W^{1,p}(U)}$$
for some appropriate constant $C$, which does not depend on $u$.

4. Inequality (2) holds for $u \in C^1(\overline{U})$. Assume now $u \in W^{1,p}(U)$. Then there exist functions $u_m \in C^\infty(\overline{U})$ converging to $u$ in $W^{1,p}(U)$ (Theorem \ref{thm:global-smooth-approximation-boundary}). According to (2) we have
$$(3) \quad \|Tu_m - Tu_l\|_{L^p(\partial U)} \le C\|u_m - u_l\|_{W^{1,p}(U)};$$
so that $\{Tu_m\}_{m=1}^\infty$ is a Cauchy sequence in $L^p(\partial U)$. We define
$$Tu := \lim_{m \to \infty} Tu_m,$$
the limit taken in $L^p(\partial U)$. According to (3) this definition does not depend on the particular choice of smooth functions approximating $u$.

Finally if $u \in W^{1,p}(U) \cap C(\bar U)$, we note that the functions $u_m \in C^\infty(\bar U)$ constructed in the proof of Theorem \ref{thm:global-smooth-approximation-boundary} converge uniformly to $u$ on $\bar U$. Hence $Tu = u|_{\partial U}$.
\end{proof}

We next examine more closely what it means for a function to have zero trace.

\begin{theorem}[Trace-zero functions in $W^{1,p}$]
\label{thm:trace-zero-characterization}
\dcref{ch5:5.5-thm-2}
\group{Sobolev Traces}
\uses{def:sobolev-space-zero-boundary,thm:trace-theorem}
Assume $U$ is bounded and $\partial U$ is $C^1$. Suppose furthermore that $u \in W^{1,p}(U)$. Then
$$(4) \quad u \in W_0^{1,p}(U) \quad \text{if and only if} \quad Tu = 0 \text{ on } \partial U.$$
\end{theorem}

\begin{proof} \textit{(Omit on first reading.)}
1. Suppose first $u \in W_0^{1,p}(U)$. Then by definition there exist functions $u_m \in C_c^\infty(U)$ such that
$$u_m \to u \quad \text{in } W^{1,p}(U).$$
As $Tu_m = 0$ on $\partial U$ ($m = 1,\ldots$) and $T : W^{1,p}(U) \to L^p(\partial U)$ is a bounded linear operator, we deduce $Tu = 0$ on $\partial U$.

2. The converse statement is more difficult. Let us assume that
$$(5) \quad Tu = 0 \quad \text{on } \partial U.$$
Using partitions of unity and flattening out $\partial U$ as usual, we may as well assume
$$(6) \quad \begin{cases} u \in W^{1,p}(\R_+^n), \quad u \text{ has compact support in } \R_+^n, \\ Tu = 0 \text{ on } \partial \R_+^n = \R^{n-1}. \end{cases}$$
Then since $Tu = 0$ on $\R^{n-1}$, there exist functions $u_m \in C^1(\R_+^n)$ such that
$$(7) \quad u_m \to u \text{ in } W^{1,p}(\R_+^n)$$
and
$$(8) \quad Tu_m = u_m|_{\R^{n-1}} \to 0 \text{ in } L^p(\R^{n-1}).$$
Now if $x' \in \R^{n-1}, x_n \ge 0$, we have
$$|u_m(x', x_n)| \le |u_m(x', 0)| + \int_0^{x_n} |u_{m,x_n}(x',t)| \, dt.$$
Thus
$$\int_{\R^{n-1}} |u_m(x', x_n)|^p \, dx' \le C \left( \int_{\R^{n-1}} |u_m(x', 0)|^p \, dx' + x_n^{p-1} \int_0^{x_n} \int_{\R^{n-1}} |Du_m(x',t)|^p \, dx' \, dt \right).$$
Letting $m \to \infty$ and recalling (7), (8), we deduce:
$$(9) \quad \int_{\R^{n-1}} |u(x', x_n)|^p \, dx' \le Cx_n^{p-1} \int_0^{x_n} \int_{\R^{n-1}} |Du|^p \, dx' \, dt$$
for a.e. $x_n > 0$.

3. Next let $\zeta \in C^\infty(\R)$ satisfy
$$\zeta \equiv 1 \text{ on } [0,1], \quad \zeta \equiv 0 \text{ on } \R - [0,2], \quad 0 \le \zeta \le 1,$$
and write
$$\begin{cases} \zeta_m(x) := \zeta(mx_n) & (x \in \R_+^n) \\ w_m := u(x)(1 - \zeta_m). \end{cases}$$
Then
$$\begin{cases} w_{m,x_n} = u_{x_n}(1 - \zeta_m) - mu\zeta' \\ D_x' w_m = D_x' u(1 - \zeta_m). \end{cases}$$
Consequently
$$\int_{\R_+^n} |Dw_m - Du|^p \, dx \le C \int_{\R_+^n} |\zeta_m|^p |Du|^p \, dx$$
$$(10) \quad + Cm^p \int_0^{2/m} \int_{\R^{n-1}} |u|^p \, dx' \, dt$$
$$=: A + B.$$
Now
$$(11) \quad A \to 0 \quad \text{as } m \to \infty,$$
since $\zeta_m \ne 0$ only if $0 \le x_n \le 2/m$. To estimate the term $B$, we utilize inequality (9):
$$(12) \quad B \le Cm^p \left(\int_0^{2/m} t^{p-1} \, dt\right) \left(\int_0^{2/m} \int_{\R^{n-1}} |Du|^p \, dx' \, dx_n\right)$$
$$\le C \int_0^{2/m} \int_{\R^{n-1}} |Du|^p \, dx' \, dx_n \to 0 \quad \text{as } m \to \infty.$$
Employing (10)--(12), we deduce $Dw_m \to Du$ in $L^p(\R_+^n)$. Since clearly $w_m \to u$ in $L^p(\R_+^n)$, we conclude
$$w_m \to u \quad \text{in } W^{1,p}(\R_+^n).$$
But $w_m = 0$ if $0 < x_n < 1/m$. We can therefore mollify the $w_m$ to produce functions $u_m \in C_c^\infty(\R_+^n)$ such that $u_m \to u$ in $W^{1,p}(\R_+^n)$. Hence $u \in W_0^{1,p}(\R_+^n)$.
\end{proof}

\section{Sobolev Inequalities}
\label{sec:sobolev-inequalities}

Our goal in this section is to discover embeddings of various Sobolev spaces into others. The crucial analytic tools here will be certain so-called ``Sobolev-type inequalities'', which we will prove below for smooth functions. These will then establish the estimates for arbitrary functions in the various relevant Sobolev spaces, since---as we saw in \cref{sec:sobolev-approximation}---smooth functions are dense.

To clarify the presentation we will consider first only the Sobolev space $W^{1,p}(U)$ and ask the following basic question: \textit{If a function $u$ belongs to $W^{1,p}(U)$, does $u$ automatically belong to certain other spaces?} The answer will be ``yes'', but which other spaces depends upon whether
$$(1) \quad 1 \le p < n,$$
$$(2) \quad p = n,$$
$$(3) \quad n < p \le \infty.$$
We study case (1) in \cref{sec:sobolev-inequalities} (§5.6.1), case (3) in §5.6.2, and the borderline case (2) only later in §5.8.1.

\subsection{Gagliardo--Nirenberg--Sobolev inequality}

For this section let us assume
$$(4) \quad 1 \le p < n$$
and first ask whether we can establish an estimate of the form
$$(5) \quad \|u\|_{L^q(\R^n)} \le C\|Du\|_{L^p(\R^n)},$$
for certain constants $C > 0$, $1 \le q < \infty$ and all functions $u \in C_c^\infty(\R^n)$. The point is that the constants $C$ and $q$ should not depend on $u$.

\textbf{Motivation.} Let us first demonstrate that if any inequality of the form (5) holds, then the number $q$ cannot be arbitrary, but must in fact have a very specific form. For this, choose a function $u \in C_c^\infty(\R^n)$, $u \ne 0$, and define for $\lambda > 0$ the rescaled function
$$u_\lambda(x) := u(\lambda x) \quad (x \in \R^n).$$
Applying (5) to $u_\lambda$, we find:
$$(6) \quad \|u_\lambda\|_{L^q(\R^n)} \le C\|Du_\lambda\|_{L^p(\R^n)}.$$
Now
$$\int_{\R^n} |u_\lambda|^q\,dx = \int_{\R^n} |u(\lambda x)|^q\,dx = \frac{1}{\lambda^n} \int_{\R^n} |u(y)|^q\,dy,$$
and
$$\int_{\R^n} |Du_\lambda|^p\,dx = \lambda^p \int_{\R^n} |Du(\lambda x)|^p\,dx = \frac{\lambda^p}{\lambda^n} \int_{\R^n} |Du(y)|^p\,dy.$$
Inserting these equalities into (6), we discover
$$\frac{1}{\lambda^{n/q}} \|u\|_{L^q(\R^n)} \le C \frac{\lambda}{\lambda^{n/p}} \|Du\|_{L^p(\R^n)},$$
and so
$$(7) \quad \|u\|_{L^q(\R^n)} \le C\lambda^{1 - n/p + n/q} \|Du\|_{L^p(\R^n)}.$$
But then if $1 - \frac{n}{p} + \frac{n}{q} \ne 0$, we can upon sending $\lambda$ to either $0$ or $\infty$ in (7) obtain a contradiction. Thus if in fact the desired inequality (5) holds, we must necessarily have $1 - \frac{n}{p} + \frac{n}{q} = 0$; so that $\frac{1}{q} = \frac{1}{p} - \frac{1}{n}$, $q = \frac{np}{n-p}$.

This observation motivates the following

\begin{definition}[Sobolev conjugate exponent]
\label{def:sobolev-conjugate-exponent}
\dcref{ch5:5.6-def-1}
\group{Sobolev Inequalities}
If $1 \le p < n$, the Sobolev conjugate of $p$ is
$$(8) \quad p^* := \frac{np}{n - p}.$$
Note that
$$(9) \quad \frac{1}{p^*} = \frac{1}{p} - \frac{1}{n}, \quad p^* > p.$$
\end{definition}

The foregoing scaling analysis shows the estimate (5) can only possibly be true for $q = p^*$. Next we prove this inequality is in fact valid.

\begin{theorem}[Gagliardo--Nirenberg--Sobolev inequality]
\label{thm:gagliardo-nirenberg-sobolev-inequality}
\dcref{ch5:5.6-thm-1}
\group{Sobolev Inequalities}
\uses{def:sobolev-conjugate-exponent}
Assume $1 \le p < n$. There exists a constant $C$, depending only on $p$ and $n$, such that
$$(10) \quad \|u\|_{L^{p^*}(\R^n)} \le C\|Du\|_{L^p(\R^n)},$$
for all $u \in C_c^1(\R^n)$.
\end{theorem}

Now we really do need $u$ to have compact support for (10) to hold, as the example $u \equiv 1$ shows. But remarkably the constant here does not depend at all upon the size of the support of $u$.

\begin{proof}
1. First assume $p = 1$.

Since $u$ has compact support, for each $i = 1,\ldots,n$ and $x \in \R^n$ we have
$$u(x) = \int_{-\infty}^{x_i} u_{x_i}(x_1,\ldots,x_{i-1},y_i,x_{i+1},\ldots,x_n)\,dy_i,$$
and so
$$|u(x)| \le \int_{-\infty}^{\infty} |Du(x_1,\ldots,y_i,\ldots,x_n)|\,dy_i \quad (i=1,\ldots,n).$$
Consequently
$$(11) \quad |u(x)|^{\frac{n}{n-1}} \le \prod_{i=1}^n \left(\int_{-\infty}^\infty |Du(x_1,\ldots,y_i,\ldots,x_n)|\,dy_i\right)^{\frac{1}{n-1}}.$$
Integrate this inequality with respect to $x_1$:
$$(12) \quad \int_{-\infty}^\infty |u|^{\frac{n}{n-1}}\,dx_1 \le \int_{-\infty}^\infty \prod_{i=1}^n \left(\int_{-\infty}^\infty |Du|\,dy_i\right)^{\frac{1}{n-1}} dx_1$$
$$= \left(\int_{-\infty}^\infty |Du|\,dy_1\right)^{\frac{1}{n-1}} \int_{-\infty}^\infty \prod_{i=2}^n \left(\int_{-\infty}^\infty |Du|\,dy_i\right)^{\frac{1}{n-1}} dx_1$$
$$\le \left(\int_{-\infty}^\infty |Du|\,dy_1\right)^{\frac{1}{n-1}} \left(\prod_{i=2}^n \int_{-\infty}^\infty \int_{-\infty}^\infty |Du|\,dx_1\,dy_i\right)^{\frac{1}{n-1}},$$
the last inequality resulting from the general Hölder inequality (§B.2).

Now integrate (12) with respect to $x_2$:
$$\int_{-\infty}^{\infty} \int_{-\infty}^{\infty} |u|^{\frac{n}{n-1}} dx_1 dx_2 \le \left(\int_{-\infty}^{\infty} \int_{-\infty}^{\infty} |Du| dx_1 dy_2\right)^{\frac{1}{n-1}} \int_{-\infty}^{\infty} \prod_{\substack{i=1 \\ i \ne 2}}^{n} I_i^{\frac{1}{n-1}} dx_2,$$
for
$$I_1 := \int_{-\infty}^{\infty} |Du| dy_1, \quad I_i := \int_{-\infty}^{\infty} \int_{-\infty}^{\infty} |Du| dx_1 dy_i \quad (i = 3, \ldots, n).$$
Applying once more the extended Hölder inequality, we find
$$\int_{-\infty}^{\infty} \int_{-\infty}^{\infty} |u|^{\frac{n}{n-1}} dx_1 dx_2 \le \left(\int_{-\infty}^{\infty} \int_{-\infty}^{\infty} |Du| dx_1 dy_2\right)^{\frac{1}{n-1}} \left(\int_{-\infty}^{\infty} \int_{-\infty}^{\infty} |Du| dy_1 dx_2\right)^{\frac{1}{n-1}}$$
$$\prod_{i=3}^{n} \left(\int_{-\infty}^{\infty} \int_{-\infty}^{\infty} \int_{-\infty}^{\infty} |Du| dx_1 dx_2 dy_i\right)^{\frac{1}{n-1}}.$$
We continue by integrating with respect to $x_3, \ldots, x_n$, eventually to find
$$(13) \quad \int_{\R^n} |u|^{\frac{n}{n-1}} dx \le \prod_{i=1}^{n} \left(\int_{-\infty}^{\infty} \cdots \int_{-\infty}^{\infty} |Du| dx_1 \ldots dy_i \ldots dx_n\right)^{\frac{1}{n-1}} = \left(\int_{\R^n} |Du| dx\right)^{\frac{n}{n-1}}.$$
This is estimate (10) for $p = 1$.

2. Consider now the case that $1 < p < n$. We apply estimate (13) to $v := |u|^\gamma$, where $\gamma > 1$ is to be selected. Then
$$(14) \quad \left(\int_{\R^n} |u|^{\frac{\gamma n}{n-1}} dx\right)^{\frac{n-1}{n}} \le \int_{\R^n} |D|u|^\gamma| dx = \gamma \int_{\R^n} |u|^{\gamma-1} |Du| dx$$
$$\le \gamma \left(\int_{\R^n} |u|^{(\gamma-1)\frac{p}{p-1}} dx\right)^{\frac{p-1}{p}} \left(\int_{\R^n} |Du|^p dx\right)^{\frac{1}{p}}.$$
We choose $\gamma$ so that $\frac{\gamma n}{n-1} = (\gamma - 1)\frac{p}{p-1}$. That is, we set
$$\gamma = \frac{p(n-1)}{n-p} > 1;$$
in which case $\frac{\gamma n}{n-1} = (\gamma - 1)\frac{p}{p-1} = \frac{np}{n-p} = p^*$. Thus, in view of (5), estimate (14) becomes
$$\left(\int_{\R^n} |u|^{p^*} dx\right)^{\frac{1}{p^*}} \le C\left(\int_{\R^n} |Du|^p dx\right)^{\frac{1}{p}}.$$
\end{proof}

\begin{theorem}[Estimates for $W^{1,p}$, $1 \le p < n$]
\label{thm:sobolev-embedding-w1p-subcritical}
\dcref{ch5:5.6-thm-2}
\group{Sobolev Inequalities}
\uses{thm:gagliardo-nirenberg-sobolev-inequality,thm:sobolev-extension-theorem,thm:local-smooth-approximation}
Let $U$ be a bounded, open subset of $\R^n$, and suppose $\partial U$ is $C^1$. Assume $1 \le p < n$, and $u \in W^{1,p}(U)$. Then $u \in L^{p^*}(U)$, with the estimate
$$(15) \quad \|u\|_{L^{p^*}(U)} \le C\|u\|_{W^{1,p}(U)},$$
the constant $C$ depending only on $p, n$, and $U$.
\end{theorem}

\begin{proof}
Since $\partial U$ is $C^1$, there exists according to Theorem \ref{thm:sobolev-extension-theorem} an extension $E u = \bar{u} \in W^{1,p}(\R^n)$, such that
$$(16) \quad \begin{cases}
\bar{u} = u \text{ in } U, \bar{u} \text{ has compact support, and} \\
\|\bar{u}\|_{W^{1,p}(\R^n)} \le C\|u\|_{W^{1,p}(U)}.
\end{cases}$$
Because $\bar{u}$ has compact support, we know from Theorem \ref{thm:local-smooth-approximation} that there exist functions $u_m \in C_c^\infty(\R^n)$ ($m = 1, 2, \ldots$) such that
$$(17) \quad u_m \to \bar{u} \text{ in } W^{1,p}(\R^n).$$
Now according to Theorem \ref{thm:gagliardo-nirenberg-sobolev-inequality}, $\|u_m - u_l\|_{L^{p^*}(\R^n)} \le C\|Du_m - Du_l\|_{L^p(\R^n)}$ for all $l, m \ge 1$. Thus
$$(18) \quad u_m \to \bar{u} \text{ in } L^{p^*}(\R^n)$$
as well. Since Theorem \ref{thm:gagliardo-nirenberg-sobolev-inequality} also implies $\|u_m\|_{L^{p^*}(\R^n)} \le C\|Du_m\|_{L^p(\R^n)}$, assertions (17) and (18) yield the bound
$$\|\bar{u}\|_{L^{p^*}(\R^n)} \le C\|D\bar{u}\|_{L^p(\R^n)}.$$
This inequality and (16) complete the proof.
\end{proof}

\begin{theorem}[Poincaré's inequality for $W_0^{1,p}$, $1 \le p < n$]
\label{thm:poincare-w01p-subcritical}
\dcref{ch5:5.6-thm-3}
\group{Sobolev Inequalities}
\uses{def:sobolev-space-zero-boundary,thm:gagliardo-nirenberg-sobolev-inequality}
Assume $U$ is a bounded, open subset of $\R^n$. Suppose $u \in W_0^{1,p}(U)$ for some $1 \le p < n$. Then we have the estimate
$$\|u\|_{L^q(U)} \le C\|Du\|_{L^p(U)}$$
for each $q \in [1, p^*]$, the constant $C$ depending only on $p, q, n$ and $U$.
\end{theorem}

\begin{remark}[Poincaré's inequality]
\label{rem:poincare-w01p-terminology}
\dcref{ch5:5.6-rem-1}
\group{Sobolev Inequalities}
\uses{thm:poincare-w01p-subcritical,thm:sobolev-embedding-w1p-subcritical}
This estimate is sometimes called \textit{Poincaré's inequality}. The difference with Theorem \ref{thm:sobolev-embedding-w1p-subcritical} is that only the gradient of $u$ appears on the righthand side of the inequality. (Other Poincaré-type inequalities will be established later, in §5.8.1.)
\end{remark}

\begin{proof}
Since $u \in W_0^{1,p}(U)$, there exist functions $u_m \in C_c^\infty(U)$ ($m = 1,2,\ldots$) converging to $u$ in $W^{1,p}(U)$. We extend each function $u_m$ to be 0 on $\R^n - \bar{U}$ and apply Theorem \ref{thm:gagliardo-nirenberg-sobolev-inequality} to discover $\|u\|_{L^{p^*}(U)} \le C\|Du\|_{L^p(U)}$. As $|U| < \infty$, we furthermore have $\|u\|_{L^q(U)} \le C\|u\|_{L^{p^*}(U)}$ if $1 \le q \le p^*$.
\end{proof}

\begin{remark}[Equivalent norm on $W_0^{1,p}$; the borderline case $p=n$]
\label{rem:borderline-case-p-equals-n}
\dcref{ch5:5.6-rem-2}
\group{Sobolev Inequalities}
\uses{thm:poincare-w01p-subcritical}
(i) In view of Theorem \ref{thm:poincare-w01p-subcritical}, on $W_0^{1,p}(U)$ the norm $\|Du\|_{L^p(U)}$ is equivalent to $\|u\|_{W^{1,p}(U)}$, if $U$ is bounded.

(ii) We next consider the case
$$p = n.$$
Owing to Theorem \ref{thm:sobolev-embedding-w1p-subcritical} and the fact that $p^* = \frac{np}{n-p} \to +\infty$ as $p \to n$, we might expect $u \in L^\infty(U)$, provided $u \in W^{1,n}(U)$. This is however false if $n > 1$: for example, if $U = B^0(0,1)$ the function $u = \log\log\left(1 + \frac{1}{|x|}\right)$ belongs to $W^{1,n}(U)$, but not to $L^\infty(U)$. We will return to this borderline situation in §5.8.1 below.
\end{remark}

\subsection{Morrey's inequality}

Now let us suppose
$$(19) \quad n < p < \infty.$$
We will show that if $u \in W^{1,p}(U)$, then $u$ is in fact Hölder continuous, after possibly being redefined on a set of measure zero.

\begin{theorem}[Morrey's inequality]
\label{thm:morreys-inequality}
\dcref{ch5:5.6-thm-4}
\group{Sobolev Inequalities}
\uses{def:holder-seminorm-sup-norm}
Assume $n < p \le \infty$. Then there exists a constant $C$, depending only on $p$ and $n$, such that
$$(20) \quad \|u\|_{C^{0,\gamma}(\R^n)} \le C\|u\|_{W^{1,p}(\R^n)}$$
for all $u \in C^1(\R^n)$, where
$$\gamma := 1 - n/p.$$
\end{theorem}

\begin{proof}
1. First choose any ball $B(x,r) \subset \R^n$.

We claim there exists a constant $C$, depending only on $n$, such that
$$(21) \quad \int_{B(x,r)} |u(y) - u(x)|\,dy \le C \int_{B(x,r)} \frac{|Du(y)|}{|y-x|^{n-1}}\,dy.$$
To prove this, fix any point $w \in \partial B(0, 1)$. Then if $0 < s < r$,
$$|u(x + sw) - u(x)| = \left|\int_0^s \frac{d}{dt} u(x + tw)\,dt\right| = \left|\int_0^s Du(x + tw) \cdot w\,dt\right| \le \int_0^s |Du(x + tw)|\,dt.$$
Hence
$$\int_{\partial B(0,1)} |u(x + sw) - u(x)|\,dS \le \int_0^s \int_{\partial B(0,1)} |Du(x + tw)|\,dS\,dt = \int_0^s \int_{\partial B(0,1)} |Du(x + tw)| \frac{t^{n-1}}{t^{n-1}}\,dS\,dt.$$
Let $y = x + tw$, so that $t = |x - y|$. Then converting from polar coordinates, we have
$$\int_{\partial B(0,1)} |u(x + sw) - u(x)|\,dS \le \int_{B(x,s)} \frac{|Du(y)|}{|x - y|^{n-1}}\,dy \le \int_{B(x,r)} \frac{|Du(y)|}{|x - y|^{n-1}}\,dy.$$
Multiply by $s^{n-1}$ and integrate from 0 to $r$ with respect to $s$:
$$\int_{B(x,r)} |u(y) - u(x)|\,dy \le \frac{r^n}{n} \int_{B(x,r)} \frac{|Du(y)|}{|x - y|^{n-1}}\,dy.$$
This is (21).

2. Now fix $x \in \R^n$. We apply inequality (21) as follows:
$$(22) \quad |u(x)| \le \int_{B(x,1)} |u(x) - u(y)|\,dy + \int_{B(x,1)} |u(y)|\,dy$$
$$\le C \int_{B(x,1)} \frac{|Du(y)|}{|x - y|^{n-1}}\,dy + C\|u\|_{L^p(B(x,1))}$$
$$\le C \left(\int_{\R^n} |Du|^p\,dy\right)^{1/p} \left(\int_{B(x,1)} \frac{dy}{|x - y|^{(n-1)\frac{p}{p-1}}}\right)^{\frac{p-1}{p}} + C\|u\|_{L^p(\R^n)}$$
$$\le C\|u\|_{W^{1,p}(\R^n)}.$$
The last estimate holds since $p > n$ implies $(n-1)\frac{p}{p-1} < n$; so that
$$\int_{B(x,1)} \frac{1}{|x - y|^{(n-1)\frac{p}{p-1}}}\,dy < \infty.$$
As $x \in \R^n$ is arbitrary, inequality (22) implies
$$(23) \quad \sup_{\R^n} |u| \le C\|u\|_{W^{1,p}(\R^n)}.$$

3. Next, choose any two points $x, y \in \R^n$ and write $r := |x - y|$. Let $W := B(x,r) \cap B(y,r)$. Then
$$(24) \quad |u(x) - u(y)| \le \int_W |u(x) - u(z)|\,dz + \int_W |u(y) - u(z)|\,dz.$$
But inequality (21) allows us to estimate
$$\int_W |u(x) - u(z)|\,dz \le C \int_{B(x,r)} |u(x) - u(z)|\,dz$$
$$(25) \quad \le C \left( \int_{B(x,r)} |Du|^p\,dz \right)^{1/p} \left( \int_{B(x,r)} \frac{dz}{|x-z|^{(n-1)p/(p-1)}} \right)^{(p-1)/p}$$
$$\le C \left( r^{n-(n-1)\frac{p}{p-1}} \right)^{\frac{p-1}{p}} \|Du\|_{L^p(\R^n)}$$
$$= Cr^{1-n/p} \|Du\|_{L^p(\R^n)}.$$
Likewise,
$$\int_W |u(y) - u(z)|\,dz \le Cr^{1-n/p} \|Du\|_{L^p(\R^n)}.$$
On substituting this estimate and (25) into (24) yields
$$|u(x) - u(y)| \le Cr^{1-n/p}\|Du\|_{L^p(\R^n)} = C|x-y|^{1-n/p}\|Du\|_{L^p(\R^n)}.$$
Thus
$$[u]_{C^{0,1-n/p}(\R^n)} = \sup_{x \ne y} \left\{ \frac{|u(x) - u(y)|}{|x-y|^{1-n/p}} \right\} \le C\|Du\|_{L^p(\R^n)}.$$
This inequality and (23) complete the proof of (20).
\end{proof}

\begin{remark}[Morrey estimate on a ball]
\label{rem:morrey-estimate-on-ball}
\dcref{ch5:5.6-rem-3}
\group{Sobolev Inequalities}
\uses{thm:morreys-inequality}
A slight variant of the proof above provides the estimate
$$|u(y) - u(x)| \le Cr^{1-n/p} \left( \int_{B(x,2r)} |Du(z)|^p\,dz \right)^{1/p}$$
for all $u \in C^1(B(x, 2r))$, $y \in B(x, r)$, $n < p < \infty$. By approximation the same bound is valid for $u \in W^{1,p}(B(x, 2r))$, $n < p < \infty$. We will use this inequality later in §5.8.2. (This estimate is in fact valid if on the right hand side we integrate over $B(x, r)$, instead of $B(x, 2r)$, but the proof is a bit trickier.)
\end{remark}

\begin{definition}[Version of a function]
\label{def:version-of-a-function}
\dcref{ch5:5.6-def-2}
\group{Sobolev Inequalities}
We say $u^*$ is a version of a given function $u$ provided
$$u = u^* \text{ a.e.}$$
\end{definition}

\begin{theorem}[Estimates for $W^{1,p}$, $n < p \le \infty$]
\label{thm:sobolev-embedding-w1p-supercritical}
\dcref{ch5:5.6-thm-5}
\group{Sobolev Inequalities}
\uses{def:version-of-a-function,thm:morreys-inequality,thm:sobolev-extension-theorem,thm:local-smooth-approximation}
Let $U$ be a bounded, open subset of $\R^n$, and suppose $\partial U$ is $C^1$. Assume $n < p \le \infty$, and $u \in W^{1,p}(U)$. Then $u$ has a version $u^* \in C^{0,\gamma}(\bar{U})$, for $\gamma = 1 - \frac{n}{p}$, with the estimate
$$\|u^*\|_{C^{0,\gamma}(\bar{U})} \le C\|u\|_{W^{1,p}(U)}.$$
The constant $C$ depends only on $p, n$ and $U$.
\end{theorem}

\begin{remark}[Identification with continuous version]
\label{rem:identify-continuous-version}
\dcref{ch5:5.6-rem-4}
\group{Sobolev Inequalities}
\uses{thm:sobolev-embedding-w1p-supercritical}
In view of Theorem \ref{thm:sobolev-embedding-w1p-supercritical}, \textit{we will henceforth always identify a function $u \in W^{1,p}(U)$ ($p > n$) with its continuous version}.
\end{remark}

\begin{proof}
Since $\partial U$ is $C^1$, there exists according to Theorem \ref{thm:sobolev-extension-theorem} an extension $E u = \bar{u} \in W^{1,p}(\R^n)$ such that
$$(26) \quad \begin{cases}
\bar{u} = u \text{ in } U, \\
\bar{u} \text{ has compact support, and} \\
\|\bar{u}\|_{W^{1,p}(\R^n)} \le C\|u\|_{W^{1,p}(U)}.
\end{cases}$$
Since $\bar{u}$ has compact support, we obtain from Theorem \ref{thm:local-smooth-approximation} the existence of functions $u_m \in C_c^\infty(\R^n)$ such that
$$(27) \quad u_m \to \bar{u} \text{ in } W^{1,p}(\R^n).$$
Now according to Theorem \ref{thm:morreys-inequality}, $\|u_m - u_l\|_{C^{0,1-n/p}(\R^n)} \le C\|u_m - u_l\|_{W^{1,p}(\R^n)}$ for all $l, m \ge 1$; whence there exists a function $u^* \in C^{0,1-n/p}(\R^n)$ such that
$$(28) \quad u_m \to u^* \text{ in } C^{0,1-n/p}(\R^n).$$
Owing to (27) and (28) we see that $u^* = u$ a.e. on $U$; so that $u^*$ is a version of $u$. Since Theorem \ref{thm:morreys-inequality} also implies $\|u_m\|_{C^{0,1-n/p}(\R^n)} \le C\|u_m\|_{W^{1,p}(\R^n)}$, assertions (27) and (28) yield:
$$\|u^*\|_{C^{0,1-n/p}(\R^n)} \le C\|\bar{u}\|_{W^{1,p}(\R^n)}.$$
This inequality and (26) complete the proof.
\end{proof}

\subsection{General Sobolev inequalities}

We can now concatenate the estimates established in §§5.6.1 and 5.6.2 to obtain more complicated (and hard-to-remember) inequalities.

\begin{theorem}[General Sobolev inequalities]
\label{thm:general-sobolev-inequalities}
\dcref{ch5:5.6-thm-6}
\group{Sobolev Inequalities}
\uses{thm:gagliardo-nirenberg-sobolev-inequality,thm:morreys-inequality}
Let $U$ be a bounded open subset of $\R^n$, with a $C^1$ boundary. Assume $u \in W^{k,p}(U)$.

(i) If
$$(29) \quad k < \frac{n}{p},$$
then $u \in L^q(U)$, where
$$\frac{1}{q} = \frac{1}{p} - \frac{k}{n}.$$
We have in addition the estimate
$$(30) \quad \|u\|_{L^q(U)} \le C\|u\|_{W^{k,p}(U)},$$
the constant $C$ depending only on $k, p, n$ and $U$.

(ii) If
$$(31) \quad k > \frac{n}{p},$$
then $u \in C^{k-[\frac{n}{p}]-1,\gamma}(\bar{U})$, where
$$\gamma = \begin{cases}
\left[\frac{n}{p}\right] + 1 - \frac{n}{p}, & \text{if } \frac{n}{p} \text{ is not an integer} \\
\text{any positive number} < 1, & \text{if } \frac{n}{p} \text{ is an integer}.
\end{cases}$$
We have in addition the estimate
$$(32) \quad \|u\|_{C^{k-[\frac{n}{p}]-1,\gamma}(\bar{U})} \le C\|u\|_{W^{k,p}(U)},$$
the constant $C$ depending only on $k, p, n, \gamma$ and $U$.
\end{theorem}

\begin{proof}
1. Assume (29). Then since $D^\alpha u \in L^p(U)$ for all $|\alpha| = k$, the Sobolev--Nirenberg--Gagliardo inequality implies
$$\|D^\beta u\|_{L^{p^*}(U)} \le C\|u\|_{W^{k,p}(U)} \quad \text{if } |\beta| = k-1,$$
and so $u \in W^{k-1,p^*}(U)$. Similarly, we find $u \in W^{k-2,p^{**}}(U)$, where $\frac{1}{p^{**}} = \frac{1}{p^*} - \frac{1}{n} = \frac{1}{p} - \frac{2}{n}$. Continuing, we eventually discover after $k$ steps that $u \in W^{0,q}(U) = L^q(U)$, for $\frac{1}{q} = \frac{1}{p} - \frac{k}{n}$. The stated estimate (30) follows from multiplying the relevant estimates at each stage of the above argument.

2. Assume now condition (31) holds, and $\frac{n}{p}$ is not an integer. Then as above we see
$$(33) \quad u \in W^{k-1,r}(U)$$
for
$$(34) \quad \frac{1}{r} = \frac{1}{p} - \frac{1}{n},$$
provided $lp < n$. We choose the integer $l$ so that
$$(35) \quad l < \frac{n}{p} < l+1;$$
that is, we set $l = \left[\frac{n}{p}\right]$. Consequently (34) and (35) imply $r = \frac{pn}{n-pl} > n$. Hence (33) and Morrey's inequality imply that $D^\alpha u \in C^{0,1-\frac{n}{r}}(\bar{U})$ for all $|\alpha| \le k - l - 1$. Observe also that $1 - \frac{n}{r} = 1 - \frac{n}{p} + l = \left[\frac{n}{p}\right] + 1 - \frac{n}{p}$. Thus
$$u \in C^{k-\left[\frac{n}{p}\right]-1,\left[\frac{n}{p}\right]+1-\frac{n}{p}}(\bar{U}),$$
and the stated estimate follows easily.

3. Finally, suppose (31) holds, with $\frac{n}{p}$ an integer. Set $l = \left[\frac{n}{p}\right]-1 = \frac{n}{p}-1$. Consequently, we have as above $u \in W^{k-l,r}(U)$ for $r = \frac{pn}{n-pl} = n$. Hence the Sobolev--Nirenberg--Gagliardo inequality shows $D^\alpha u \in L^q(U)$ for all $n \le q < \infty$ and all $|\alpha| \le k-l-1 = k-\left[\frac{n}{p}\right]$. Therefore Morrey's inequality further implies $D^\alpha u \in C^{0,1-\frac{n}{q}}(\bar U)$ for all $n < q < \infty$ and all $|\alpha| \le k-\left[\frac{n}{p}\right]-1$. Consequently $u \in C^{k-\left[\frac{n}{p}\right]-1,\gamma}(\bar U)$ for each $0 < \gamma < 1$. As before, the stated estimate follows as well.
\end{proof}

\begin{remark}[Fourier-transform proof]
\label{rem:general-sobolev-fourier-remark}
\dcref{ch5:5.6-rem-5}
\group{Sobolev Inequalities}
\uses{thm:general-sobolev-inequalities}
Various general Sobolev-type inequalities can also be proved using the Fourier transform: see Problem 18.
\end{remark}

\section{Compactness}
\label{sec:sobolev-compactness}

We have seen in \cref{sec:sobolev-inequalities} that the Gagliardo--Nirenberg--Sobolev inequality implies the embedding of $W^{1,p}(U)$ into $L^{p^*}(U)$ for $1 \le p < n$, $p^* = \frac{pn}{n-p}$. We will now demonstrate that $W^{1,p}(U)$ is in fact \textit{compactly} embedded in $L^q(U)$ for $1 \le q < p^*$. This compactness will be fundamental for our applications of linear and nonlinear functional analysis to PDE in Chapters 6-9.

\begin{definition}[Compact embedding]
\label{def:compact-embedding}
\dcref{ch5:5.7-def-1}
\group{Compactness}
Let $X$ and $Y$ be Banach spaces, $X \subset Y$. We say that $X$ is compactly embedded in $Y$, written
$$X \CC Y,$$
provided
\begin{itemize}
\item[(i)] $\|x\|_Y \le C\|x\|_X$ ($x \in X$) for some constant $C$,
\end{itemize}
and

(ii) each bounded sequence in $X$ is precompact in $Y$.
\end{definition}

\begin{theorem}[Rellich--Kondrachov Compactness Theorem]
\label{thm:rellich-kondrachov}
\dcref{ch5:5.7-thm-1}
\group{Compactness}
\uses{def:compact-embedding,thm:sobolev-embedding-w1p-subcritical,thm:sobolev-extension-theorem,thm:local-smooth-approximation}
Assume $U$ is a bounded open subset of $\R^n$, and $\partial U$ is $C^1$. Suppose $1 \le p < n$. Then
$$W^{1,p}(U) \CC L^q(U)$$
for each $1 \le q < p^*$.
\end{theorem}

\begin{proof}
1. Fix $1 \le q < p^*$ and note that since $U$ is bounded, Theorem \ref{thm:sobolev-embedding-w1p-subcritical} implies
$$W^{1,p}(U) \subset L^q(U), \quad \|u\|_{L^q(U)} \le C\|u\|_{W^{1,p}(U)}.$$
It remains therefore to show that if $\{u_m\}_{m=1}^\infty$ is a bounded sequence in $W^{1,p}(U)$, there exists a subsequence $\{u_{m_j}\}_{j=1}^\infty$ which converges in $L^q(U)$.

2. In view of the Extension Theorem from \cref{sec:sobolev-extensions} we may with no loss of generality assume that $U = \R^n$ and the functions $\{u_m\}_{m=1}^\infty$ all have compact support in some bounded open set $V \subset \R^n$. We also may assume
$$(1) \quad \sup_m \|u_m\|_{W^{1,p}(V)} < \infty.$$

3. Let us first study the smoothed functions
$$u_m^\varepsilon := \eta_\varepsilon * u_m \quad (\varepsilon > 0, m = 1,2,\ldots),$$
$\eta_\varepsilon$ denoting the usual mollifier. We may suppose the functions $\{u_m^\varepsilon\}_{m=1}^\infty$ all have support in $V$ as well.

4. We first claim
$$(2) \quad u_m^\varepsilon \to u_m \text{ in } L^q(V) \text{ as } \varepsilon \to 0, \text{ uniformly in } m.$$
To prove this, we first note that if $u_m$ is smooth, then
$$u_m^\varepsilon(x) - u_m(x) = \int_{B(0,1)} \eta(y)(u_m(x-\varepsilon y) - u_m(x)) \, dy$$
$$= \int_{B(0,1)} \eta(y) \int_0^1 \frac{d}{dt}(u_m(x - \varepsilon ty)) \, dt\, dy$$
$$= -\varepsilon \int_{B(0,1)} \eta(y) \int_0^1 Du_m(x - \varepsilon ty) \cdot y \, dt\, dy.$$
Thus
$$\int_V |u_m^\varepsilon(x) - u_m(x)| dx \le \varepsilon \int_{B(0,1)} \eta(y) \int_0^1 \int_V |Du_m(x - \varepsilon ty)| dx dt dy$$
$$\le \varepsilon \int_V |Du_m(z)| dz.$$
By approximation this estimate holds if $u_m \in W^{1,p}(V)$. Hence
$$\|u_m^\varepsilon - u_m\|_{L^1(V)} \le \varepsilon\|Du_m\|_{L^1(V)} \le \varepsilon C\|Du_m\|_{L^p(V)},$$
the latter inequality holding since $V$ is bounded. Owing to (1) we thereby discover
$$(3) \quad u_m^\varepsilon \to u_m \text{ in } L^1(V), \text{ uniformly in } m.$$
But then since $1 \le q < p^*$, we see using the interpolation inequality for $L^p$-norms (§B.2) that
$$\|u_m^\varepsilon - u_m\|_{L^q(V)} \le \|u_m^\varepsilon - u_m\|_{L^1(V)}^\theta \|u_m^\varepsilon - u_m\|_{L^{p^*}(V)}^{1-\theta},$$
where $\frac{1}{q} = \theta + \frac{(1-\theta)}{p^*}$, $0 < \theta < 1$. Consequently (1) and the Gagliardo--Nirenberg--Sobolev inequality imply
$$\|u_m^\varepsilon - u_m\|_{L^q(V)} \le C\|u_m^\varepsilon - u_m\|_{L^1(V)}^\theta;$$
whence assertion (2) follows from (3).

5. Next we claim
$$(4) \quad \text{for each fixed } \varepsilon > 0, \text{ the sequence } \{u_m^\varepsilon\}_{m=1}^\infty \text{ is uniformly bounded and equicontinuous.}$$
Indeed, if $x \in \R^n$, then
$$|u_m^\varepsilon(x)| \le \int_{B(x,\varepsilon)} \eta_\varepsilon(x-y)|u_m(y)| dy$$
$$\le \|\eta_\varepsilon\|_{L^\infty(\R^n)}\|u_m\|_{L^1(V)} \le \frac{C}{\varepsilon^n} < \infty$$
for $m = 1, 2, \ldots$ Similarly
$$|Du_m^\varepsilon(x)| \le \int_{B(x,\varepsilon)} |D\eta_\varepsilon(x-y)||u_m(y)| dy$$
$$\le \|D\eta_\varepsilon\|_{L^\infty(\R^n)}\|u_m\|_{L^1(V)} \le \frac{C}{\varepsilon^{n+1}} < \infty,$$
for $m = 1,\ldots$. Assertion (4) follows from these two estimates.

6. Now fix $\delta > 0$. We will show there exists a subsequence $\{u_{m_j}\}_{j=1}^\infty \subset \{u_m\}_{m=1}^\infty$ such that
$$(5) \quad \limsup_{j,k\to\infty} \|u_{m_j} - u_{m_k}\|_{L^q(V)} \le \delta.$$
To see this, let us first employ assertion (2) to select $\varepsilon > 0$ so small that
$$(6) \quad \|u_m^\varepsilon - u_m\|_{L^q(V)} \le \delta/2$$
for $m = 1, 2, \ldots$.

We now observe that since the functions $\{u_m\}_{m=1}^\infty$, and thus the functions $\{u_m^\varepsilon\}_{m=1}^\infty$, have support in some fixed bounded set $V \subset \R^n$, we may utilize (4) and the Arzelà--Ascoli compactness criterion, §C.7, to obtain a subsequence $\{u_{m_j}^\varepsilon\}_{j=1}^\infty \subset \{u_m^\varepsilon\}_{m=1}^\infty$ which converges \textit{uniformly} on $V$. In particular therefore
$$(7) \quad \limsup_{j,k\to\infty} \|u_{m_j}^\varepsilon - u_{m_k}^\varepsilon\|_{L^q(V)} = 0.$$
But then (6) and (7) imply
$$\limsup_{j,k\to\infty} \|u_{m_j} - u_{m_k}\|_{L^q(V)} \le \delta,$$
and so (5) is proved.

7. We next employ assertion (5) with $\delta = 1, \frac{1}{2}, \frac{1}{3}, \ldots$ and use a standard diagonal argument to extract a subsequence $\{u_{m_l}\}_{l=1}^\infty \subset \{u_m\}_{m=1}^\infty$ satisfying
$$\limsup_{l,k\to\infty} \|u_{m_l} - u_{m_k}\|_{L^q(V)} = 0.$$
\end{proof}

\begin{remark}[$W^{1,p}(U) \CC L^p(U)$ for all $p$]
\label{rem:compact-embedding-lp-all-p}
\dcref{ch5:5.7-rem-1}
\group{Compactness}
\uses{thm:rellich-kondrachov,thm:sobolev-embedding-w1p-subcritical,thm:morreys-inequality}
Observe that since $p^* > p$ and $p^* \to \infty$ as $p \to n$, we have in particular
$$W^{1,p}(U) \CC L^p(U)$$
for all $1 \le p \le \infty$. (Observe that if $n < p \le \infty$, this follows from Morrey's inequality and the Arzelà--Ascoli compactness criterion.) Note also that
$$W_0^{1,p}(U) \CC L^p(U),$$
even if we do not assume $\partial U$ to be $C^1$.
\end{remark}

\section{Additional Topics}
\label{sec:sobolev-additional-topics}

\subsection{Poincaré's inequalities}

We now illustrate how the compactness assertion in \cref{sec:sobolev-compactness} can be used to generate new inequalities.

\textbf{Notation.} $(u)_U = \fint_U u\,dy =$ average of $u$ over $U$.

\begin{theorem}[Poincaré's inequality]
\label{thm:poincares-inequality-connected}
\dcref{ch5:5.8-thm-1}
\group{Poincaré Inequalities}
\uses{thm:rellich-kondrachov,def:sobolev-space-wkp}
Let $U$ be a bounded, connected, open subset of $\R^n$, with a $C^1$ boundary $\partial U$. Assume $1 \le p \le \infty$. Then there exists a constant $C$, depending only on $n,p$ and $U$, such that
$$(1) \quad \|u - (u)_U\|_{L^p(U)} \le C\|Du\|_{L^p(U)}$$
for each function $u \in W^{1,p}(U)$.
\end{theorem}

The significance of (1) is that only the gradient of $u$ appears on the right hand side.

\begin{proof}
We argue by contradiction. Were the stated estimate false, there would exist for each integer $k = 1,\ldots$ a function $u_k \in W^{1,p}(U)$ satisfying
$$(2) \quad \|u_k - (u_k)_U\|_{L^p(U)} > k\|Du_k\|_{L^p(U)}.$$
We renormalize by defining
$$(3) \quad v_k := \frac{u_k - (u_k)_U}{\|u_k - (u_k)_U\|_{L^p(U)}} \quad (k = 1,\ldots).$$
Then
$$(v_k)_U = 0, \quad \|v_k\|_{L^p(U)} = 1;$$
and (2) implies
$$(4) \quad \|Dv_k\|_{L^p(U)} < \frac{1}{k} \quad (k = 1,2,\ldots).$$
In particular the functions $\{v_k\}_{k=1}^\infty$ are bounded in $W^{1,p}(U)$.

In view of the Remark after the Rellich--Kondrachov Theorem in \cref{sec:sobolev-compactness}, there exists a subsequence $\{v_{k_j}\}_{j=1}^\infty \subset \{v_k\}_{k=1}^\infty$ and a function $v \in L^p(U)$ such that
$$(5) \quad v_{k_j} \to v \quad \text{in } L^p(U).$$
From (3) it follows that
$$(6) \quad (v)_U = 0, \quad \|v\|_{L^p(U)} = 1.$$
On the other hand, (4) implies for each $i = 1, \ldots, n$ and $\phi \in C_c^\infty(U)$ that
$$\int_U v\phi_{x_i}\, dx = \lim_{k_j \to \infty} \int_U v_{k_j}\phi_{x_i}\, dx = -\lim_{k_j \to \infty} \int_U v_{k_j,x_i}\phi\, dx = 0.$$
Consequently $v \in W^{1,p}(U)$, with $Dv = 0$ a.e. Thus $v$ is constant, since $U$ is connected (see Problem 10). However this conclusion is at variance with (6): since $v$ is constant and $(v)_U = 0$, we must have $v = 0$; in which case $\|v\|_{L^p(U)} = 0$. This contradiction establishes estimate (1).
\end{proof}

A particularly important special case follows.

\textbf{Notation.} $(u)_{x,r} = \fint_{B(x,r)} u\,dy =$ average of $u$ over the ball $B(x,r)$.

\begin{theorem}[Poincaré's inequality for a ball]
\label{thm:poincares-inequality-ball}
\dcref{ch5:5.8-thm-2}
\group{Poincaré Inequalities}
\uses{thm:poincares-inequality-connected}
Assume $1 \le p \le \infty$. Then there exists a constant $C$, depending only on $n$ and $p$, such that
$$(7) \quad \|u - (u)_{x,r}\|_{L^p(B(x,r))} \le Cr\|Du\|_{L^p(B(x,r))}$$
for each ball $B(x,r) \subset \R^n$ and each function $u \in W^{1,p}(B^0(x,r))$.
\end{theorem}

\begin{proof}
1. The case $U = B^0(0,1)$ follows from Theorem \ref{thm:poincares-inequality-connected}. In general, if $u \in W^{1,p}(B^0(x,r))$ write
$$v(y) := u(x + ry) \quad (y \in B(0,1)).$$
Then $v \in W^{1,p}(B^0(0,1))$, and we have
$$\|v - (v)_{0,1}\|_{L^p(B(0,1))} \le C\|Dv\|_{L^p(B(0,1))}.$$
Changing variables, we recover estimate (7).
\end{proof}

\begin{remark}[Bounded mean oscillation]
\label{rem:bmo-remark}
\dcref{ch5:5.8-rem-1}
\group{Poincaré Inequalities}
\uses{thm:poincares-inequality-ball}
Assume $u \in W^{1,n}(\R^n) \cap L^1(\R^n)$, and let $B(x,r)$ be any ball. Then Theorem \ref{thm:poincares-inequality-ball} with $p = 1$ implies
$$\int_{B(x,r)} |u - (u)_{x,r}|\, dy \le Cr\int_{B(x,r)} |Du|\, dy$$
$$\le Cr \left(\int_{B(x,r)} |Du|^n\, dy\right)^{1/n} \le C\left(\int_{\R^n} |Du|^n\, dy\right)^{1/n}.$$
Thus $u \in BMO(\R^n)$, the space of functions of \textit{bounded mean oscillation} in $\R^n$, with the seminorm
$$[u]_{BMO(\R^n)} := \sup_{B(x,r) \subset \R^n} \left\{\int_{B(x,r)} |u - (u)_{x,r}|\, dy\right\}.$$
See Stein \cite{Stein1993} for the theory of $BMO$.
\end{remark}

\subsection{Difference quotients}

When we later apply Sobolev space theory to PDE, we will be forced to study difference quotient approximations to weak derivatives. Following is the relevant theory, which the reader may wish to postpone studying until the need arises, in §6.3.

\textbf{a. Difference quotients and $W^{1,p}$.}

Assume $u : U \to \R$ is a locally summable function, and $V \CC U$.

\begin{definition}[Difference quotients]
\label{def:difference-quotients}
\dcref{ch5:5.8-def-1}
\group{Difference Quotients}
(i) The $i^{\text{th}}$-difference quotient of size $h$ is
$$D_i^h u(x) = \frac{u(x + he_i) - u(x)}{h} \quad (i = 1, \ldots, n)$$
for $x \in V$ and $h \in \R$, $0 < |h| < \dist(V, \partial U)$.

(ii) $D^h u := (D_1^h u, \ldots, D_n^h u)$.
\end{definition}

\begin{theorem}[Difference quotients and weak derivatives]
\label{thm:difference-quotients-weak-derivatives}
\dcref{ch5:5.8-thm-3}
\group{Difference Quotients}
\uses{def:difference-quotients,def:sobolev-space-wkp}
(i) Suppose $1 \le p < \infty$ and $u \in W^{1,p}(U)$. Then for each $V \CC U$
$$(8) \quad \|D^h u\|_{L^p(V)} \le C\|Du\|_{L^p(U)}$$
for some constant $C$ and all $0 < |h| < \frac{1}{2}\dist(V, \partial U)$.

(ii) Assume $1 < p < \infty$, $u \in L^p(V)$, and there exists a constant $C$ such that
$$(9) \quad \|D^h u\|_{L^p(V)} \le C$$
for all $0 < |h| < \frac{1}{2}\dist(V, \partial U)$. Then
$$u \in W^{1,p}(V), \quad \text{with } \|Du\|_{L^p(V)} \le C.$$
\end{theorem}

\begin{remark}[Assertion (ii) fails for $p=1$]
\label{rem:difference-quotient-p1-fails}
\dcref{ch5:5.8-rem-2}
\group{Difference Quotients}
\uses{thm:difference-quotients-weak-derivatives}
Assertion (ii) is false for $p = 1$ (Problem 11).
\end{remark}

\begin{proof}
1. Assume $1 \le p < \infty$, and temporarily suppose $u$ is smooth. Then for each $x \in V$, $i = 1, \ldots, n$, and $0 < |h| < \frac{1}{2}\dist(V, \partial U)$, we have
$$u(x + he_i) - u(x) = \int_0^1 u_{x_i}(x + the_i) \, dt \cdot h;$$
so that
$$|u(x + he_i) - u(x)| \le h \int_0^1 |Du(x + the_i)| \, dt.$$
Consequently
$$\int_V |D^h u|^p\,dx \le C \sum_{i=1}^n \int_V \int_0^1 |Du(x + the_i)|^p\,dt\,dx$$
$$= C \sum_{i=1}^n \int_0^1 \int_V |Du(x + the_i)|^p\,dx\,dt.$$
Thus
$$\int_V |D^h u|^p\,dx \le C \int_U |Du|^p\,dx.$$
This estimate holds should $u$ be smooth, and thus is valid by approximation for arbitrary $u \in W^{1,p}(U)$.

2. Now suppose estimate (9) holds for all $0 < |h| < \frac{1}{2}\dist(V, \partial U)$ and some constant $C$. Choose $i = 1, \ldots, n$, $\phi \in C_c^\infty(V)$, and note for small enough $h$ that
$$\int_V u(x) \left[\frac{\phi(x + he_i) - \phi(x)}{h}\right] dx = -\int_V \left[\frac{u(x) - u(x - he_i)}{h}\right] \phi(x)\,dx;$$
that is,
$$(10) \quad \int_V u(D_i^h \phi)\,dx = -\int_V (D_i^{-h} u)\phi\,dx.$$
This is the ``integration-by-parts'' formula for difference quotients. Estimate (9) implies
$$\sup_h \|D_i^{-h} u\|_{L^p(V)} < \infty;$$
and therefore, since $1 < p < \infty$, there exists a function $v_i \in L^p(V)$ and a subsequence $h_k \to 0$ such that
$$(11) \quad D_i^{-h_k} u \rightharpoonup v_i \quad \text{weakly in } L^p(V).$$
(See §D.4 for weak convergence.) But then
$$\int_V u\phi_{x_i}\,dx = \int_U u\phi_{x_i}\,dx = \lim_{h_k\to 0} \int_U u D_i^{h_k}\phi\,dx$$
$$= -\lim_{h_k\to 0} \int_V D_i^{-h_k}u\,\phi\,dx$$
$$= -\int_V v_i\phi\,dx = -\int_U v_i\phi\,dx.$$
Thus $v_i = u_{x_i}$ in the weak sense ($i = 1, \ldots, n$), and so $Du \in L^p(V)$. As $u \in L^p(V)$, we deduce therefore that $u \in W^{1,p}(V)$.
\end{proof}

\begin{remark}[Variants of Theorem \ref{thm:difference-quotients-weak-derivatives} for half-balls]
\label{rem:difference-quotient-half-ball-remark}
\dcref{ch5:5.8-rem-3}
\group{Difference Quotients}
\uses{thm:difference-quotients-weak-derivatives}
Variants of Theorem \ref{thm:difference-quotients-weak-derivatives} can be valid even if it is not true that $V \CC U$. For example if $U$ is the open half-ball $B^0(0,1) \cap \{x_n > 0\}$, $V = B^0(0,1/2) \cap \{x_n > 0\}$, we have the bound $\int_V |D_i^h u|^p\,dx \le \int_U |u_{x_i}|^p\,dx$ for $i = 1, \ldots, n-1$. The proof is similar to that just given.

We will need this comment in Chapter 6, §6.3.2.
\end{remark}

\textbf{b. Lipschitz functions and $W^{1,\infty}$.}

\begin{theorem}[Characterization of $W^{1,\infty}$]
\label{thm:characterization-w1infty-lipschitz}
\dcref{ch5:5.8-thm-4}
\group{Difference Quotients}
\uses{def:sobolev-space-wkp,thm:sobolev-extension-theorem}
Let $U$ be open and bounded, with $\partial U$ of class $C^1$. Then $u : U \to \R$ is Lipschitz continuous if and only if $u \in W^{1,\infty}(U)$.
\end{theorem}

\begin{proof}
1. First assume $U = \R^n$ and $u$ has compact support.

Suppose $u \in W^{1,\infty}(\R^n)$. Then $u^\varepsilon := \eta_\varepsilon * u$, where $\eta_\varepsilon$ is the usual mollifier, is smooth and satisfies
$$\begin{cases}
u^\varepsilon \to u \text{ uniformly as } \varepsilon \to 0, \\
\|Du^\varepsilon\|_{L^\infty(\R^n)} \le \|Du\|_{L^\infty(\R^n)}.
\end{cases}$$
Choose any two points $x, y \in \R^n$, $x \ne y$. We have
$$u^\varepsilon(x) - u^\varepsilon(y) = \int_0^1 \frac{d}{dt} u^\varepsilon(tx + (1-t)y) \, dt$$
$$= \int_0^1 Du^\varepsilon(tx + (1-t)y) \, dt \cdot (x - y);$$
and so
$$|u^\varepsilon(x) - u^\varepsilon(y)| \le \|Du^\varepsilon\|_{L^\infty(\R^n)}|x - y| \le \|Du\|_{L^\infty(\R^n)}|x - y|.$$
We let $\varepsilon \to 0$ to discover
$$|u(x) - u(y)| \le \|Du\|_{L^\infty(\R^n)}|x - y|.$$
Hence $u$ is Lipschitz continuous.

2. On the other hand assume now $u$ is Lipschitz continuous; we must prove that $u$ has essentially bounded weak first derivative. Since $u$ is Lipschitz, we see
$$\|D_i^{-h}u\|_{L^\infty(\R^n)} \le \text{Lip}(u),$$
and thus there exists a function $v_i \in L^\infty(\R^n)$ and a subsequence $h_k \to 0$ such that
$$(12) \quad D_i^{-h_k}u \to v_i \quad \text{weakly in } L^2_{\text{loc}}(\R^n).$$
Consequently
$$\int_{\R^n} u\phi_{x_i}\,dx = \lim_{h_k \to 0} \int_{\R^n} uD_i^{h_k}\phi\,dx$$
$$= -\lim_{h_k \to 0} \int_{\R^n} D_i^{-h_k}u\,\phi\,dx = -\int_{\R^n} v_i\phi\,dx$$
by (12). The above equality holds for all $\phi \in C_c^\infty(\R^n)$, and so $v_i = u_{x_i}$ in the weak sense ($i = 1, \ldots, n$). Consequently $u \in W^{1,\infty}(\R^n)$.

3. In the general case that $U$ is bounded, with $\partial U$ of class $C^1$, we as usual extend $u$ to $Eu = \bar{u}$ and apply the above argument.
\end{proof}

\begin{remark}[Local Lipschitz characterization; failure for $1\le p<\infty$]
\label{rem:local-lipschitz-w1infty-loc-remark}
\dcref{ch5:5.8-rem-4}
\group{Difference Quotients}
\uses{thm:characterization-w1infty-lipschitz,thm:sobolev-embedding-w1p-supercritical}
The argument above adapts easily to prove that for any open set $U$, $u \in W^{1,\infty}_{\text{loc}}(U)$ if and only if $u$ is locally Lipschitz continuous in $U$. There is no corresponding characterization of the spaces $W^{1,p}$ for $1 \le p < \infty$. If $n < p < \infty$, then each function $u \in W^{1,p}$ belongs to $C^{0,1-n/p}$, but on the other hand a function Hölder continuous with exponent less than one need not belong to any Sobolev space $W^{1,p}$.
\end{remark}

\subsection{Differentiability a.e.}

Next we examine more closely the connections between weak partial derivatives and partial derivatives in the usual calculus sense.

\begin{definition}[Differentiability at a point]
\label{def:differentiability-at-a-point}
\dcref{ch5:5.8-def-2}
\group{Differentiability of Sobolev Functions}
A function $u: U \to \R$ is differentiable at $x \in U$ if there exists $a \in \R^n$ such that
$$(13) \quad u(y) = u(x) + a \cdot (y - x) + o(|y - x|) \quad \text{as } y \to x.$$
In other words,
$$\lim_{y \to x} \frac{|u(y) - u(x) - a \cdot (y - x)|}{|y - x|} = 0.$$
\end{definition}

\textbf{Notation.} It is easy to check that $a$, if it exists, is unique. We henceforth write
$$Du(x)$$
for $a$ and call $Du$ the gradient of $u$.

To be sure that this notation is consistent, we need to study the relationships between the various notions of derivatives:

\begin{theorem}[Differentiability almost everywhere]
\label{thm:differentiability-almost-everywhere}
\dcref{ch5:5.8-thm-5}
\group{Differentiability of Sobolev Functions}
\uses{def:differentiability-at-a-point,thm:morreys-inequality}
Assume $u \in W^{1,p}_{\text{loc}}(U)$ for some $n < p \le \infty$. Then $u$ is differentiable a.e. in $U$, and its gradient equals its weak gradient a.e.
\end{theorem}

Recall we always identify $u$ with its continuous version.

\begin{proof}
1. Assume first $n < p < \infty$. From the Remark after the proof of Theorem \ref{thm:morreys-inequality}, we recall Morrey's estimate
$$(14) \quad |v(y) - v(x)| \le Cr^{1-n/p} \left(\int_{B(x,2r)} |Dv(z)|^p\,dz\right)^{1/p} \quad (y \in B(x,r)),$$
valid for any $C^1$ function $v$ and thus, by approximation, for any $v \in W^{1,p}$.

2. Choose $u \in W^{1,p}_{\text{loc}}(U)$. Now for a.e. $x \in U$, a version of Lebesgue's Differentiation Theorem (§E.4) implies
$$(15) \quad \int_{B(x,r)} |Du(x) - Du(z)|^p\,dz \to 0$$
as $r \to 0$, $Du$ denoting as usual the weak derivative of $u$. Fix any such point $x$ and set
$$v(y) := u(y) - u(x) - Du(x) \cdot (y - x)$$
in estimate (14), where
$$(16) \quad r = |x - y|.$$
We find
$$|u(y) - u(x) - Du(x) \cdot (y - x)|$$
$$\le Cr^{1-n/p} \left(\int_{B(x,2r)} |Du(x) - Du(z)|^p\,dz\right)^{1/p}$$
$$\le Cr \left(\int_{B(x,2r)} |Du(x) - Du(z)|^p\,dz\right)^{1/p}$$
$$= o(r) \text{ by (15)}$$
$$= o(|x - y|) \text{ by (16)}.$$
Thus $u$ is differentiable at $x$, and its gradient equals its weak gradient at $x$.

3. In case $p = \infty$, we note $W^{1,\infty}_{\text{loc}}(U) \subset W^{1,p}_{\text{loc}}(U)$ for all $1 \le p < \infty$, and apply the reasoning above.
\end{proof}

Finally, in view of Theorem \ref{thm:differentiability-almost-everywhere}, we obtain

\begin{theorem}[Rademacher's Theorem]
\label{thm:rademachers-theorem}
\dcref{ch5:5.8-thm-6}
\group{Differentiability of Sobolev Functions}
\uses{thm:differentiability-almost-everywhere,thm:characterization-w1infty-lipschitz}
Let $u$ be locally Lipschitz continuous in $U$. Then $u$ is differentiable almost everywhere in $U$.
\end{theorem}

\subsection{Fourier transform methods}

Next we employ the Fourier transform (§4.3) to give an alternate characterization of the spaces $H^k(\R^n)$. For this section all functions are complex-valued.

\begin{theorem}[Characterization of $H^k$ by Fourier transform]
\label{thm:characterization-hk-fourier-transform}
\dcref{ch5:5.8-thm-7}
\group{Fourier Characterization of Sobolev Spaces}
\uses{rem:hk-notation-identification}
Let $k$ be a nonnegative integer.

(i) A function $u \in L^2(\R^n)$ belongs to $H^k(\R^n)$ if and only if
$$(1 + |y|^k)\hat{u} \in L^2(\R^n).$$

(ii) In addition, there exists a positive constant $C$ such that
$$\frac{1}{C}\|u\|_{H^k(\R^n)} \le \|(1 + |y|^k)\hat{u}\|_{L^2(\R^n)} \le C\|u\|_{H^k(\R^n)}$$
for each $u \in H^k(\R^n)$.
\end{theorem}

\begin{proof}
1. Assume first $u \in H^k(\R^n)$. Then for each multiindex $|\alpha| \le k$, we have $D^\alpha u \in L^2(\R^n)$. Now if $u \in C^k$ has compact support, we have
$$\widehat{D^\alpha u} = (iy)^\alpha \hat{u}$$
according to Theorem 2 in §4.3.1. Approximating by smooth functions we deduce this formula provided $u \in H^k(\R^n)$. Thus $(iy)^\alpha \hat{u} \in L^2(\R^n)$ for each $|\alpha| \le k$. In particular choosing $\alpha = (k,0,\ldots,0), (0,k,\ldots,0), \ldots, (0,\ldots,k)$, we deduce
$$\int_{\R^n} |y|^{2k}|\hat{u}|^2\,dy \le C \int_{\R^n} |D^k u|^2\,dx < \infty.$$
Thus
$$\int_{\R^n} (1 + |y|^k)^2 |\hat{u}|^2\,dy \le C\|u\|_{H^k(\R^n)},$$
and so $(1 + |y|^k)\hat{u} \in L^2(\R^n)$.

2. Suppose conversely $(1 + |y|^k)\hat{u} \in L^2(\R^n)$ and $|\alpha| \le k$. Then
$$\|(iy)^\alpha \hat{u}\|_{L^2(\R^n)} \le \int_{\R^n} |y|^{2|\alpha|} |\hat{u}|^2\,dy \le C\|(1 + |y|^k)\hat{u}\|_{L^2(\R^n)}.$$
Set
$$u_\alpha := ((iy)^\alpha \hat{u})^\vee.$$
Then for each $\phi \in C_c^\infty(\R^n)$
$$\int_{\R^n} (D^\alpha \phi)\, u\,dx = \int_{\R^n} (\widehat{D^\alpha \phi})\,\hat{u}\,dy = \int_{\R^n} (iy)^\alpha \hat{\phi}\, \hat{u}\,dy$$
$$= (-1)^{|\alpha|} \int_{\R^n} \phi\, u_\alpha\,dx.$$
Thus $u_\alpha = D^\alpha u$ in the weak sense and, by (20), $D^\alpha u \in L^2(\R^n)$. Hence $u \in H^k(U)$, as required.
\end{proof}

It is sometimes useful to define also \textit{fractional} Sobolev spaces.

\begin{definition}[Fractional Sobolev space $H^s(\R^n)$]
\label{def:fractional-sobolev-space}
\dcref{ch5:5.8-def-3}
\group{Fourier Characterization of Sobolev Spaces}
\uses{thm:characterization-hk-fourier-transform}
Assume $0 < s < \infty$ and $u \in L^2(\R^n)$. Then $u \in H^s(\R^n)$ if $(1 + |y|^s)\hat u \in L^2(\R^n)$. For noninteger $s$, we set
$$\|u\|_{H^s(\R^n)} := \|(1 + |y|^s)\hat u\|_{L^2(\R^n)}.$$
\end{definition}

\section{Other Spaces of Functions}
\label{sec:other-spaces-of-functions}

\subsection{The space $H^{-1}$}

As we will see later in our systematic study in Chapters 6 and 7 of linear elliptic, parabolic and hyperbolic PDE, it is important to have an explicit characterization of the dual space of $H_0^1$. (See Appendix D for definitions.)

\begin{definition}[The dual space $H^{-1}(U)$]
\label{def:dual-space-h-minus-1}
\dcref{ch5:5.9-def-1}
\group{Dual Sobolev Spaces}
\uses{def:sobolev-space-zero-boundary}
We denote by $H^{-1}(U)$ the dual space to $H_0^1(U)$.
\end{definition}

In other words $f$ belongs to $H^{-1}(U)$ provided $f$ is a bounded linear functional on $H_0^1(U)$.

\textbf{Notation.} We will write $\langle\ ,\ \rangle$ to denote the pairing between $H^{-1}(U)$ and $H_0^1(U)$.

\begin{definition}[Norm on $H^{-1}(U)$]
\label{def:norm-on-h-minus-1}
\dcref{ch5:5.9-def-2}
\group{Dual Sobolev Spaces}
\uses{def:dual-space-h-minus-1}
If $f \in H^{-1}(U)$, we define the norm
$$\|f\|_{H^{-1}(U)} = \sup\left\{\langle f,u\rangle \mid u \in H_0^1(U), \|u\|_{H_0^1(U)} \le 1\right\}.$$
\end{definition}

\begin{theorem}[Characterization of $H^{-1}$]
\label{thm:characterization-h-minus-1}
\dcref{ch5:5.9-thm-1}
\group{Dual Sobolev Spaces}
\uses{def:norm-on-h-minus-1}
(i) Assume $f \in H^{-1}(U)$. Then there exist functions $f^0, f^1, \ldots, f^n$ in $L^2(U)$ such that
$$(1) \quad \langle f,v \rangle = \int_U f^0 v + \sum_{i=1}^n f^i v_{x_i}\,dx \quad (v \in H_0^1(U)).$$

(ii) Furthermore,
$$\|f\|_{H^{-1}(U)} = \inf\left\{\left(\int_U \sum_{i=0}^n |f^i|^2\,dx\right)^{1/2} \;\middle|\; f \text{ satisfies (1) for } f^0,\ldots,f^n \in L^2(U)\right\}.$$
\end{theorem}

\textbf{Notation.} We write $f = f^0 - \sum_{i=1}^n f_{x_i}^i$ whenever (1) holds.

\begin{proof}
1. Given $u, v \in H_0^1(U)$, we define the inner product $(u,v) := \int_U Du \cdot Dv + uv\,dx$. Let $f \in H^{-1}(U)$. We apply the Riesz Representation Theorem (§D.3) to deduce the existence of a unique function $w \in H_0^1(U)$ satisfying $B[u,v] = (f,v)$ for all $v \in H_0^1(U)$; that is,
$$\int_U Du \cdot Dv + uv\,dx = (f,v) \tag{2}$$
for each $v \in H_0^1(U)$. This establishes (1) for
$$\begin{cases} f^0 = u \\ f^i = u_{x_i} \quad (i = 1, \ldots, n). \end{cases} \tag{3}$$

2. Assume now $f \in H^{-1}(U)$,
$$\langle f, v \rangle = \int_U g^0 v + \sum_{i=1}^n g^i v_{x_i}\,dx \tag{4}$$
for $g^0, g^1, \ldots, g^n \in L^2(U)$. Setting $v = u$ in (2) and using (4), we deduce
$$\int_U |Du|^2 + |u|^2\,dx \le \int_U \sum_{i=0}^n |g^i|^2\,dx.$$
Thus (3) implies
$$\int_U \sum_{i=0}^n |f^i|^2\,dx \le \int_U \sum_{i=0}^n |g^i|^2\,dx. \tag{5}$$

3. From (1) it follows that
$$|\langle f, v \rangle| \le \left(\int_U \sum_{i=0}^n |f^i|^2\,dx\right)^{1/2}$$
if $\|v\|_{H_0^1(U)} \le 1$. Consequently
$$\|f\|_{H^{-1}(U)} \le \left(\int_U \sum_{i=0}^n |f^i|^2\,dx\right)^{1/2}.$$
Setting $v = \frac{u}{\|u\|_{H_0^1(U)}}$ in (2), we deduce that in fact
$$\|f\|_{H^{-1}(U)} = \left(\int_U \sum_{i=0}^n |f^i|^2\,dx\right)^{1/2}. \tag{6}$$
Assertion (ii) follows now from (4)--(6).
\end{proof}

\subsection{Spaces involving time}

We study next some other sorts of Sobolev spaces, these comprising functions mapping time into Banach spaces. These will prove essential in our constructions of weak solutions to linear parabolic and hyperbolic PDE in Chapter 7 and to nonlinear parabolic PDE in Chapter 9.

Let $X$ denote a real Banach space, with norm $\|\cdot\|$. The reader should first of all read §5.5 about measure and integration theory for mappings taking values in $X$.

\begin{definition}[$L^p(0,T;X)$]
\label{def:lp-time-banach-space}
\dcref{ch5:5.9-def-3}
\group{Bochner Spaces}
The space
$$L^p(0, T; X)$$
consists of all measurable functions $\mathbf{u} : [0, T] \to X$ with
\begin{itemize}
\item[(i)] $\|\mathbf{u}\|_{L^p(0,T;X)} := \left(\int_0^T \|\mathbf{u}(t)\|^p\,dt\right)^{1/p} < \infty$
\end{itemize}
for $1 \le p < \infty$, and
\begin{itemize}
\item[(ii)] $\|\mathbf{u}\|_{L^\infty(0,T;X)} := \text{ess sup}_{0 \le t \le T} \|\mathbf{u}(t)\| < \infty$.
\end{itemize}
\end{definition}

\begin{definition}[$C([0,T];X)$]
\label{def:continuous-time-banach-space}
\dcref{ch5:5.9-def-4}
\group{Bochner Spaces}
The space
$$C([0, T]; X)$$
comprises all continuous functions $\mathbf{u} : [0, T] \to X$ with
$$\|\mathbf{u}\|_{C([0,T];X)} := \max_{0 \le t \le T} \|\mathbf{u}(t)\| < \infty.$$
\end{definition}

\begin{definition}[Weak derivative for Banach-space-valued functions]
\label{def:weak-derivative-banach-valued}
\dcref{ch5:5.9-def-5}
\group{Bochner Spaces}
Let $\mathbf{u} \in L^1(0, T; X)$. We say $\mathbf{v} \in L^1(0, T; X)$ is the weak derivative of $\mathbf{u}$, written
$$\mathbf{u}' = \mathbf{v},$$
provided
$$\int_0^T \phi'(t)\mathbf{u}(t)\,dt = -\int_0^T \phi(t)\mathbf{v}(t)\,dt$$
for all scalar test functions $\phi \in C_c^\infty(0, T)$.
\end{definition}

\begin{definition}[$W^{1,p}(0,T;X)$ and $H^1(0,T;X)$]
\label{def:sobolev-space-time-banach}
\dcref{ch5:5.9-def-6}
\group{Bochner Spaces}
\uses{def:lp-time-banach-space,def:weak-derivative-banach-valued}
(i) The Sobolev space
$$W^{1,p}(0,T;X)$$
consists of all functions $u \in L^p(0,T;X)$ such that $u'$ exists in the weak sense and belongs to $L^p(0,T;X)$. Furthermore,
$$\|u\|_{W^{1,p}(0,T;X)} := \begin{cases}
\left(\int_0^T \|u(t)\|^p + \|u'(t)\|^p \, dt\right)^{1/p} & (1 \le p < \infty) \\
\text{ess sup}_{0 \le t \le T}(\|u(t)\| + \|u'(t)\|) & (p = \infty).
\end{cases}$$

(ii) We write $H^1(0,T;X) = W^{1,2}(0,T;X)$.
\end{definition}

\begin{theorem}[Calculus in an abstract space]
\label{thm:calculus-abstract-space}
\dcref{ch5:5.9-thm-2}
\group{Bochner Spaces}
\uses{def:sobolev-space-time-banach,thm:local-smooth-approximation}
Let $u \in W^{1,p}(0,T;X)$ for some $1 \le p \le \infty$. Then

(i) $u \in C([0,T];X)$ (after possibly being redefined on a set of measure zero), and

(ii) $u(t) = u(s) + \int_s^t u'(\tau) \, d\tau$ for all $0 \le s \le t \le T$.

(iii) Furthermore, we have the estimate
$$(7) \quad \max_{0 \le t \le T} \|u(t)\| \le C\|u\|_{W^{1,p}(0,T;X)},$$
the constant $C$ depending only on $T$.
\end{theorem}

\begin{proof}
1. Extend $u$ to be 0 on $(-\infty, 0)$ and $(T, \infty)$, and then set $u^\varepsilon = \eta_\varepsilon * u$, $\eta_\varepsilon$ denoting the usual mollifier on $\R^1$. We check as in the proof of Theorem \ref{thm:local-smooth-approximation} that $(u^\varepsilon)' = \eta_\varepsilon * u'$ on $(\varepsilon, T - \varepsilon)$.

Then as $\varepsilon \to 0$,
$$(8) \quad \begin{cases}
u^\varepsilon \to u & \text{in } L^p(0,T;X), \\
(u^\varepsilon)' \to u' & \text{in } L^p(0,T;X).
\end{cases}$$
Fixing $0 < s < t < T$, we compute
$$u^\varepsilon(t) = u^\varepsilon(s) + \int_s^t (u^\varepsilon)'(\tau) \, d\tau.$$
Thus
$$(9) \quad u(t) = u(s) + \int_s^t u'(\tau) \, d\tau$$
for a.e. $0 < s < t < T$, according to (8). As the mapping $t \mapsto \int_0^t u'(\tau) \, d\tau$ is continuous, assertions (i), (ii) follow.

2. Estimate (7) follows easily from (9).
\end{proof}

The next two propositions concern what happens when $u$ and $u'$ lie in different spaces.

\begin{theorem}[More calculus]
\label{thm:more-calculus-h1-hminus1}
\dcref{ch5:5.9-thm-3}
\group{Bochner Spaces}
\uses{def:sobolev-space-time-banach,def:dual-space-h-minus-1}
Suppose $u \in L^2(0,T;H_0^1(U))$, with $u' \in L^2(0,T;H^{-1}(U))$.

(i) Then
$$u \in C([0,T];L^2(U))$$
(after possibly being redefined on a set of measure zero).

(ii) The mapping
$$t \mapsto \|u(t)\|^2_{L^2(U)}$$
is absolutely continuous, with
$$\frac{d}{dt}\|u(t)\|^2_{L^2(U)} = 2\langle u'(t), u(t)\rangle$$
for a.e. $0 \le t \le T$.

(iii) Furthermore, we have the estimate
$$(10) \quad \max_{0 \le t \le T} \|u(t)\|_{L^2(U)} \le C\left(\|u\|_{L^2(0,T;H_0^1(U))} + \|u'\|_{L^2(0,T;H^{-1}(U))}\right),$$
the constant $C$ depending only on $T$.
\end{theorem}

\begin{proof}
1. Extend $u$ to the larger interval $[-\sigma, T + \sigma]$ for $\sigma > 0$, and define the regularizations $u^\varepsilon = \eta_\varepsilon * u$, as in the earlier proof. Then for $\varepsilon, \delta > 0$,
$$\frac{d}{dt}\|u^\varepsilon(t) - u^\delta(t)\|^2_{L^2(U)} = 2((u^\varepsilon)'(t) - (u^\delta)'(t), u^\varepsilon(t) - u^\delta(t))_{L^2(U)}.$$
Thus
$$(11) \quad \|u^\varepsilon(t) - u^\delta(t)\|^2_{L^2(U)} = \|u^\varepsilon(s) - u^\delta(s)\|^2_{L^2(U)} + 2 \int_s^t ((u^\varepsilon)'(\tau) - (u^\delta)'(\tau), u^\varepsilon(\tau) - u^\delta(\tau)) \, d\tau$$
for all $0 \le s, t \le T$. Fix any point $s \in (0,T)$ for which
$$u^\varepsilon(s) \to u(s) \quad \text{in } L^2(U).$$
Consequently (11) implies
$$\limsup_{\varepsilon,\delta \to 0} \sup_{0 \le t \le T} \|u^\varepsilon(t) - u^\delta(t)\|^2_{L^2(U)} \le \lim_{\varepsilon,\delta \to 0} \int_0^T \|(u^\varepsilon)'(\tau) - (u^\delta)'(\tau)\|^2_{H^{-1}(U)}$$
$$+ \|u^\varepsilon(\tau) - u^\delta(\tau)\|^2_{H_0^1(U)} \, d\tau = 0.$$
Thus the smoothed functions $\{u^\varepsilon\}_{0<\varepsilon<1}$ converge in $C([0,T];L^2(U))$ to a limit $\mathbf{v} \in C([0,T];L^2(U))$. Since we also know $u^\varepsilon(t) \to u(t)$ for a.e. $t$, we deduce $u = \mathbf{v}$ a.e.

2. We similarly have
$$\|u^\varepsilon(t)\|^2_{L^2(U)} = \|u^\varepsilon(s)\|^2_{L^2(U)} + 2 \int_s^t ((u^\varepsilon)'(\tau), u^\varepsilon(\tau)) \, d\tau,$$
and so, identifying $u$ with $\mathbf{v}$ above,
$$(12) \quad \|u(t)\|^2_{L^2(U)} = \|u(s)\|^2_{L^2(U)} + 2 \int_s^t \langle u'(\tau), u(\tau) \rangle \, d\tau$$
for all $0 \le s, t \le T$.

3. To obtain (10), we integrate (12) with respect to $s$, recall the inequality $|\langle u',u\rangle| \le \|u'\|_{H^{-1}(U)}\|u\|_{H^1_0(U)}$, and make some simple estimates.
\end{proof}

For use later in the regularity theory for second-order parabolic and hyperbolic equations in Chapter 7, we will also need this extension of Theorem \ref{thm:more-calculus-h1-hminus1}.

\begin{theorem}[Mappings into better spaces]
\label{thm:mappings-into-better-spaces}
\dcref{ch5:5.9-thm-4}
\group{Bochner Spaces}
\uses{def:sobolev-space-time-banach,thm:sobolev-extension-theorem,thm:difference-quotients-weak-derivatives,thm:more-calculus-h1-hminus1}
Assume that $U$ is open, bounded, and $\partial U$ is smooth. Take $m$ to be a nonnegative integer.

Suppose $\mathbf{u} \in L^2(0,T;H^{m+2}(U))$, with $\mathbf{u}' \in L^2(0,T;H^m(U))$.

(i) Then
$$\mathbf{u} \in C([0,T];H^{m+1}(U))$$
(after possibly being redefined on a set of measure zero).

(ii) Furthermore, we have the estimate
$$(13) \quad \max_{0 \le t \le T} \|\mathbf{u}(t)\|_{H^{m+1}(U)} \le C\bigl(\|\mathbf{u}\|_{L^2(0,T;H^{m+2}(U))} + \|\mathbf{u}'\|_{L^2(0,T;H^m(U))}\bigr),$$
the constant $C$ depending only on $T$, $U$, and $m$.
\end{theorem}

\begin{proof}
1. Suppose first that $m=0$, in which case
$$\mathbf{u} \in L^2(0,T;H^2(U)), \quad \mathbf{u}' \in L^2(0,T;L^2(U)).$$
We select a bounded open set $V \supset\supset U$, and then construct a corresponding extension $\bar{\mathbf{u}} = E\mathbf{u}$, as in \cref{sec:sobolev-extensions}. In view of estimate (10) from that section, we see
$$\bar{\mathbf{u}} \in L^2(0,T;H^2(V)),$$
and
$$(14) \quad \|\bar{\mathbf{u}}\|_{L^2(0,T;H^2(V))} \le C\|\mathbf{u}\|_{L^2(0,T;H^2(U))},$$
for an appropriate constant $C$. In addition, $\bar{\mathbf{u}}' \in L^2(0,T;L^2(V))$, with the estimate
$$(15) \quad \|\bar{\mathbf{u}}'\|_{L^2(0,T;L^2(V))} \le C\|\mathbf{u}'\|_{L^2(0,T;L^2(U))}.$$
This follows if we consider difference quotients in the $t$-variable, remember the methods in §5.8.2, and observe also that $E$ is a bounded linear operator from $L^2(U)$ into $L^2(V)$.

2. Assume for the moment that $\bar{\mathbf{u}}$ is smooth. We then compute
$$\left|\frac{d}{dt}\left(\int_V |D\bar{\mathbf{u}}|^2\,dx\right)\right| = 2\left|\int_V D\bar{\mathbf{u}} \cdot D\bar{\mathbf{u}}'\,dx\right| = 2\left|\int_V \Delta\bar{\mathbf{u}}\, \bar{\mathbf{u}}'\,dx\right|$$
$$\le C\left(\|\bar{\mathbf{u}}\|^2_{H^2(V)} + \|\bar{\mathbf{u}}'\|^2_{L^2(V)}\right).$$
There is no boundary term when we integrate by parts, since the extension $\bar{\mathbf{u}} = E\mathbf{u}$ has compact support within $V$. Integrating and recalling (14), (15), it follows that
$$(16) \quad \max_{0 \le t \le T} \|\mathbf{u}(t)\|_{H^1(U)} \le C\left(\|\mathbf{u}\|_{L^2(0,T;H^2(U))} + \|\mathbf{u}'\|_{L^2(0,T;L^2(U))}\right).$$
We obtain the same estimate if $\mathbf{u}$ is not smooth, upon approximating by $\mathbf{u}^\varepsilon := \eta_\varepsilon * \mathbf{u}$, as before. As in the previous proofs, it also follows that $\mathbf{u} \in C([0,T];H^1(U))$.

3. In the general case that $m \ge 1$, we let $\alpha$ be a multiindex of order $|\alpha| \le m$, and set $\mathbf{v} := D^\alpha \mathbf{u}$. Then
$$\mathbf{v} \in L^2(0,T;H^2(U)), \quad \mathbf{v}' \in L^2(0,T;L^2(U)).$$
We apply estimate (16), with $\mathbf{v}$ replacing $\mathbf{u}$, and sum over all indices $|\alpha| \le m$, to derive (13).
\end{proof}

\section{Problems}
\label{sec:sobolev-problems}

In these exercises $U$ always denotes an open subset of $\R^n$.

%ORIGINAL-NUMBER: 1
\begin{exercise}
Suppose $k \in \{0,1,\ldots\}$, $0 < \gamma \le 1$. Prove $C^{k,\gamma}(\bar U)$ is a Banach space.
\end{exercise}

%ORIGINAL-NUMBER: 2
\begin{exercise}
Let $U$, $V$ be open sets, with $V \CC U$. Show there exists a smooth function $\zeta$ such that $\zeta \equiv 1$ on $V$, $\zeta = 0$ near $\partial U$. (Hint: Take $V \CC W \CC U$ and mollify $\chi_W$.)
\end{exercise}

%ORIGINAL-NUMBER: 3
\begin{exercise}
Assume $0 < \beta < \gamma \le 1$. Prove the interpolation inequality
$$\|u\|_{C^{0,\gamma}(U)} \le \|u\|^{\frac{1-\gamma}{1-\beta}}_{C^{0,\beta}(U)} \|u\|^{\frac{\gamma-\beta}{1-\beta}}_{C^{0,1}(U)}.$$
\end{exercise}

%ORIGINAL-NUMBER: 4
\begin{exercise}
Assume $U$ is bounded and $U \CC \bigcup_{i=1}^N V_i$. Show there exist $C^\infty$ functions $\zeta_i$ ($i = 1, \ldots, N$) such that
$$\begin{cases}
0 \le \zeta_i \le 1, \text{ spt}\,\zeta_i \subset V_i \quad (i = 1, \ldots, N) \\
\sum_{i=1}^N \zeta_i = 1 \text{ on } U.
\end{cases}$$
The functions $\{\zeta_i\}_{i=1}^N$ form a \textit{partition of unity}.
\end{exercise}

%ORIGINAL-NUMBER: 5
\begin{exercise}
Prove that if $n = 1$ and $u \in W^{1,p}(0,1)$ for some $1 \le p < \infty$, then $u$ is equal a.e. to an absolutely continuous function, and $u'$ (which exists a.e.) belongs to $L^p(0,1)$.
\end{exercise}

%ORIGINAL-NUMBER: 6
\begin{exercise}
Prove directly that if $u \in W^{1,p}(0,1)$ for some $1 < p < \infty$, then
$$|u(x) - u(y)| \le |x - y|^{1-\frac{1}{p}} \left(\int_0^1 |u'|^p\,dt\right)^{1/p} \quad \text{for a.e. } x, y \in [0,1].$$
\end{exercise}

%ORIGINAL-NUMBER: 7
\begin{exercise}
Denote by $U$ the open square $\{x \in \R^2 \mid |x_1| < 1, |x_2| < 1\}$. Define
$$u(x) = \begin{cases}
1 - x_1 & \text{if } x_1 > 0, |x_2| < x_1 \\
1 + x_1 & \text{if } x_1 < 0, |x_2| < -x_1 \\
1 - x_2 & \text{if } x_2 > 0, |x_1| < x_2 \\
1 + x_2 & \text{if } x_2 < 0, |x_1| < -x_2.
\end{cases}$$
For which $1 \le p \le \infty$ does $u$ belong to $W^{1,p}(U)$?
\end{exercise}

%ORIGINAL-NUMBER: 8
\begin{exercise}
Integrate by parts to prove the interpolation inequality:
$$\int_U |Du|^2\,dx \le C \left(\int_U u^2\,dx\right)^{1/2} \left(\int_U |D^2u|^2\,dx\right)^{1/2}$$
for all $u \in C_c^\infty(U)$. By approximation, prove this inequality if $u \in H^2(U) \cap H^1_0(U)$.
\end{exercise}

%ORIGINAL-NUMBER: 9
\begin{exercise}
Integrate by parts to prove:
$$\int_U |Du|^p\,dx \le C \left(\int_U |u|^p\,dx\right)^{1/2} \left(\int_U |D^2u|^p\,dx\right)^{1/2}$$
for $2 \le p < \infty$ and all $u \in W^{2,p}(U) \cap W^{1,p}_0(U)$. (Hint: $\int_U |Du|^p\,dx = \sum_{i=1}^n \int_U u_{x_i} u_{x_i} |Du|^{p-2}\,dx.$)
\end{exercise}

%ORIGINAL-NUMBER: 10
\begin{exercise}
Suppose $U$ is connected and $u \in W^{1,p}(U)$ satisfies
$$Du = 0 \text{ a.e. in } U.$$
Prove $u$ is constant a.e. in $U$.
\end{exercise}

%ORIGINAL-NUMBER: 11
\begin{exercise}
Show by example that if we have $\|D^h u\|_{L^1(V)} < C$ for all $0 < |h| < \frac{1}{2}\dist(V, \partial U)$, it does not necessarily follow that $u \in W^{1,1}(V)$.
\end{exercise}

%ORIGINAL-NUMBER: 12
\begin{exercise}
Give an example of an open set $U \subset \R^n$ and a function $u \in W^{1,\infty}(U)$, such that $u$ is \textit{not} Lipschitz continuous on $U$. (Hint: Take $U$ to be the open unit disk in $\R^2$, with a slit removed.)
\end{exercise}

%ORIGINAL-NUMBER: 13
\begin{exercise}
Verify that if $n > 1$, the unbounded function $u = \log \log\left(1 + \frac{1}{|x|}\right)$ belongs to $W^{1,n}(U)$, for $U = B^0(0,1)$.
\end{exercise}

%ORIGINAL-NUMBER: 14
\begin{exercise}
Let $U$ be bounded, with a $C^1$ boundary. Show that a ``typical'' function $u \in L^p(U)$ ($1 \le p < \infty$) does not have a trace on $\partial U$. More precisely, prove there does not exist a bounded linear operator
$$T : L^p(U) \to L^p(\partial U)$$
such that $Tu = u|_{\partial U}$ whenever $u \in C(\bar{U}) \cap L^p(U)$.
\end{exercise}

%ORIGINAL-NUMBER: 15
\begin{exercise}
Fix $\alpha > 0$ and let $U = B^0(0,1)$. Show there exists a constant $C$, depending only on $n$ and $\alpha$, such that
$$\int_U u^2\,dx \le C \int_U |Du|^2\,dx,$$
provided
$$|\{x \in U \mid u(x) = 0\}| \ge \alpha, \quad u \in H^1(U).$$
\end{exercise}

%ORIGINAL-NUMBER: 16
\begin{exercise}
Assume $F : \R \to \R$ is $C^1$, with $F'$ bounded. Suppose $U$ is bounded and $u \in W^{1,p}(U)$ for some $1 < p < \infty$. Show
$$v := F(u) \in W^{1,p}(U) \text{ and } v_{x_i} = F'(u) u_{x_i} \quad (i = 1, \ldots, n).$$
\end{exercise}

%ORIGINAL-NUMBER: 17
\begin{exercise}
Assume $1 < p < \infty$, and $U$ is bounded.

(i) Prove that if $u \in W^{1,p}(U)$, then $|u| \in W^{1,p}(U)$.

(ii) Prove $u \in W^{1,p}(U)$ implies $u^+, u^- \in W^{1,p}(U)$, and
$$Du^+ = \begin{cases} Du & \text{a.e. on } \{u > 0\} \\ 0 & \text{a.e. on } \{u \le 0\} \end{cases}$$
$$Du^- = \begin{cases} 0 & \text{a.e. on } \{u \ge 0\} \\ -Du & \text{a.e. on } \{u < 0\}. \end{cases}$$
(Hint: $u^+ = \lim_{\varepsilon \to 0} F_\varepsilon(u)$, for
$$F_\varepsilon(z) := \begin{cases} (z^2 + \varepsilon^2)^{1/2} - \varepsilon & \text{if } z \ge 0 \\ 0 & \text{if } z < 0. \end{cases}$$
)

(iii) Prove that if $u \in W^{1,p}(U)$, then
$$Du = 0 \text{ a.e. on the set } \{u = 0\}.$$
\end{exercise}

%ORIGINAL-NUMBER: 18
\begin{exercise}
Use the Fourier transform to prove that if $u \in H^s(\R^n)$ for $s > n/2$, then $u \in L^\infty(\R^n)$, with the bound
$$\|u\|_{L^\infty(\R^n)} \le C\|u\|_{H^s(\R^n)},$$
the constant $C$ depending only on $s$ and $n$.
\end{exercise}

\section{References}
\label{sec:sobolev-references}

\textbf{Sections 5.2--8} See Gilbarg--Trudinger \cite{GilbargTrudinger1983}, Lieb--Loss \cite{LiebLoss1997}, Ziemer \cite{Ziemer1989} and Evans--Gariepy \cite{EvansGariepy1992} for more on Sobolev spaces.

\textbf{Section 5.5} W. Schlag showed me the proof of Theorem \ref{thm:trace-zero-characterization}.

\textbf{Section 5.6} J. Ralston suggested an improvement in the proof of Theorem \ref{thm:morreys-inequality}.

\textbf{Section 5.9} See Temam \cite{Temam1977}, pp. 248-273.
